

%%%%%%%%%%%%%%%%%%%%%%%%%%%%%%%%%%%%%%%%
% diploma.tex
%
% Diplomsko delo — UL Fakulteta za računalništvo in informatiko
% Visokošolski strokovni študijski program prve stopnje
% Računalništvo in informatika
%
% Avtor: Dilan Mužič
%
% Predloga: vzorec_dip_Seminar.tex (Gašper Fijavž, Zoran Bosnić,
%           Damjan Cvetan, Franc Solina) — struktura ohranjena,
%           vsebina prilagojena temi naloge.
%
% ─────────────────────────────────────────────────────────────────
% NAVODILA ZA UPORABO
%
% 1. Ta datoteka vsebuje SAMO naslove poglavij in podpoglavij.
%    Besedilo napiši sam pod ustrezen ukaz \section / \subsection.
% 2. Poleg te datoteke potrebuješ še:
%       literatura.bib          — seznam virov
%       pdfa-1b.xmp             — metapodatki za PDF/A
%       sRGBIEC1966-2.1.icm     — barvni profil
%    Vse tri dobiš skupaj s predlogo FRI.
% 3. Prevajaj s pdflatex → bibtex → pdflatex → pdflatex.
% 4. Pred oddajo prevedi z opcijo draft (vrstica 60), da vidiš
%    predolge vrstice.
%
% Vsak stavek začni v NOVI vrstici izvorne datoteke — iskanje
% in popravljanje bosta veliko hitrejša.
% ─────────────────────────────────────────────────────────────────

\documentclass[a4paper, 12pt]{book}
\usepackage{float}
%\documentclass[a4paper, 12pt, draft]{book}  % preveri predolge vrstice!

\usepackage[utf8x]{inputenc} % slovenske črke v UTF-8
\usepackage[slovene,english]{babel} % slovenski delilni vzorci
\usepackage[pdftex]{graphicx} % vlaganje slik
\graphicspath{{slike/}} % vse slike so v mapi slike/
\usepackage{fancyhdr} % glave strani
\usepackage{amssymb} % dodatni simboli
\usepackage{amsmath} % enačbe, eqref
\usepackage[hyphens]{url}
\usepackage{booktabs} % \toprule, \midrule, \bottomrule
\usepackage{listings}
\usepackage{xcolor}
\lstset{
  basicstyle=\ttfamily\small,
  columns=fullflexible,
  breaklines=true,
  showstringspaces=false,
  frame=none,
  xleftmargin=0.6em,
  literate={č}{{\v c}}1 {š}{{\v s}}1 {ž}{{\v z}}1
           {Č}{{\v C}}1 {Š}{{\v S}}1 {Ž}{{\v Z}}1,
}

\emergencystretch=1.2em % raje malo raztegne presledke kot da besedilo uide iz roba
\usepackage{placeins} % for \FloatBarrier
\usepackage{colortbl}
\usepackage{longtable}
\usepackage{ccicons} % ikone licence Creative Commons
\usepackage{rotating} % obrnjena glava v tabeli usmerjanja
\newcommand{\rot}[1]{\rotatebox{90}{\small #1}} % tabele, ki se lahko raztezajo čez več strani
\definecolor{crta}{gray}{0.55}
\renewcommand{\topfraction}{0.9}
\renewcommand{\bottomfraction}{0.5}
\renewcommand{\textfraction}{0.07}
\renewcommand{\floatpagefraction}{0.8}
\setcounter{topnumber}{2}
\setcounter{bottomnumber}{1}
\setcounter{totalnumber}{3}

\usepackage[pdftex, colorlinks=true,
						citecolor=black, filecolor=black,
						linkcolor=black, urlcolor=black,
						pagebackref=false,
						pdfproducer={LaTeX}, pdfcreator={LaTeX}, hidelinks]{hyperref}

\usepackage{color}
\usepackage[numbers]{natbib}

%%%%%%%%%%%%%%%%%%%%%%%%%%%%%%%%%%%%%%%%
%	PODATKI O DIPLOMI
%%%%%%%%%%%%%%%%%%%%%%%%%%%%%%%%%%%%%%%%
% Če se odločiš za drug naslov, ga spremeni SAMO tukaj — ukazi
% \ttitle in \ttitleEn se uporabijo na naslovnici, v povzetku,
% v abstractu in v metapodatkih PDF.
%
% Alternativni naslovi (izberi enega, ostale zakomentiraj):
%   Večprehodna analiza argumentacije video debat z velikimi jezikovnimi modeli
%   Hibridni pristop k ocenjevanju debat: nominalne ocene jezikovnega modela in deterministični izračun
%   Spletna aplikacija za transkripcijo, preverjanje dejstev in analizo argumentacije video debat

% POZOR: slovenski in angleški naslov morata biti prevoda istega naslova —
% na ŠIS je prijavljen slovenski, angleški mora ustrezati njemu.
\newcommand{\ttitle}{Sistem za avtomatizirano analizo argumentacije video razprav z uporabo velikih jezikovnih modelov}
\newcommand{\ttitleEn}{A system for automated argumentation analysis of video debates using large language models}
\newcommand{\tsubject}{\ttitle}
\newcommand{\tsubjectEn}{\ttitleEn}
\newcommand{\tauthor}{Dilan Mužič}
\newcommand{\temail}{dm43010@student.uni-lj.si}
\newcommand{\tmentor}{izr. prof. dr. Slavko Žitnik}
\newcommand{\tyear}{2026}
\newcommand{\tkeywords}{analiza argumentacije, veliki jezikovni modeli, preverjanje dejstev, obdelava naravnega jezika, transkripcija govora}
\newcommand{\tkeywordsEn}{argumentation analysis, large language models, fact-checking, natural language processing, speech transcription}

%%%%%%%%%%%%%%%%%%%%%%%%%%%%%%%%%%%%%%%%
%	HYPERREF SETUP
%%%%%%%%%%%%%%%%%%%%%%%%%%%%%%%%%%%%%%%%
\hypersetup{pdftitle={\ttitle}}
\hypersetup{pdfsubject=\ttitleEn}
\hypersetup{pdfauthor={\tauthor, \temail}}
\hypersetup{pdfkeywords=\tkeywordsEn}

%%%%%%%%%%%%%%%%%%%%%%%%%%%%%%%%%%%%%%%%
% postavitev strani
%%%%%%%%%%%%%%%%%%%%%%%%%%%%%%%%%%%%%%%%

\addtolength{\marginparwidth}{-20pt} % robovi za tisk
\addtolength{\oddsidemargin}{40pt}
\addtolength{\evensidemargin}{-40pt}

\renewcommand{\baselinestretch}{1.3} % razmik med vrsticami
\setlength{\headheight}{15pt}
\renewcommand{\chaptermark}[1]%
{\markboth{\MakeUppercase{\thechapter.\ #1}}{}} \renewcommand{\sectionmark}[1]%
{\markright{\MakeUppercase{\thesection.\ #1}}} \renewcommand{\headrulewidth}{0.5pt} \renewcommand{\footrulewidth}{0pt}
\fancyhf{}
\fancyhead[LE,RO]{\sl \thepage}
\fancyhead[RE]{\sc \tauthor}
\fancyhead[LO]{\sc Diplomska naloga}

\newcommand{\BibTeX}{{\sc Bib}\TeX}

%%%%%%%%%%%%%%%%%%%%%%%%%%%%%%%%%%%%%%%%
% naslovi
%%%%%%%%%%%%%%%%%%%%%%%%%%%%%%%%%%%%%%%%

\newcommand{\autfont}{\Large}
\newcommand{\titfont}{\LARGE\bf}
\newcommand{\clearemptydoublepage}{\newpage{\pagestyle{empty}\cleardoublepage}}

\setcounter{tocdepth}{1}

%%%%%%%%%%%%%%%%%%%%%%%%%%%%%%%%%%%%%%%%
% konstrukti
%%%%%%%%%%%%%%%%%%%%%%%%%%%%%%%%%%%%%%%%
\newtheorem{izrek}{Izrek}[chapter]
\newtheorem{trditev}{Trditev}[izrek]
\newenvironment{dokaz}{\emph{Dokaz.}\ }{\hspace{\fill}{$\Box$}}

%%%%%%%%%%%%%%%%%%%%%%%%%%%%%%%%%%%%%%%%%%%%%%%%%%%%%%%%%%%%%%%%%%%%%%%%%%%%%%%
%% PDF-A
%%%%%%%%%%%%%%%%%%%%%%%%%%%%%%%%%%%%%%%%%%%%%%%%%%%%%%%%%%%%%%%%%%%%%%%%%%%%%%%

\def\Title{\ttitle}
\def\Author{\tauthor, \temail}
\def\Subject{\ttitleEn}
\def\Keywords{\tkeywordsEn}

%%%%%%%%%%%%%%%%%%%%%%%%%%%%%%%%%%%%%%%%
% \convertDate converts D:20080419103507+02'00' to 2008-04-19T10:35:07+02:00
%%%%%%%%%%%%%%%%%%%%%%%%%%%%%%%%%%%%%%%%
\def\convertDate{%
    \getYear
}

{\catcode`\D=12
 \gdef\getYear D:#1#2#3#4{\edef\xYear{#1#2#3#4}\getMonth}
}
\def\getMonth#1#2{\edef\xMonth{#1#2}\getDay}
\def\getDay#1#2{\edef\xDay{#1#2}\getHour}
\def\getHour#1#2{\edef\xHour{#1#2}\getMin}
\def\getMin#1#2{\edef\xMin{#1#2}\getSec}
\def\getSec#1#2{\edef\xSec{#1#2}\getTZh}
\def\getTZh #1#2{\edef\xTZh{#1#2}\getTZm}
\def\getTZm#1#2{%
    \edef\xTZm{#1#2}%
    \edef\convDate{\xYear-\xMonth-\xDay T\xHour:\xMin:\xSec+\xTZh:\xTZm}%
}

\expandafter\convertDate\pdfcreationdate

%%%%%%%%%%%%%%%%%%%%%%%%%%%%%%%%%%%%%%%%
% get pdftex version string
%%%%%%%%%%%%%%%%%%%%%%%%%%%%%%%%%%%%%%%%
\newcount\countA
\countA=\pdftexversion
\advance \countA by -100
\def\pdftexVersionStr{pdfTeX-1.\the\countA.\pdftexrevision}

%%%%%%%%%%%%%%%%%%%%%%%%%%%%%%%%%%%%%%%%
% XMP data
%%%%%%%%%%%%%%%%%%%%%%%%%%%%%%%%%%%%%%%%
\usepackage{xmpincl}
\includexmp{pdfa-1b}

%%%%%%%%%%%%%%%%%%%%%%%%%%%%%%%%%%%%%%%%
% pdfInfo
%%%%%%%%%%%%%%%%%%%%%%%%%%%%%%%%%%%%%%%%
\pdfinfo{%
    /Title (\ttitle)
    /Author (\tauthor)
    /Subject (\ttitleEn)
    /Keywords (\tkeywordsEn)
    /ModDate (\pdfcreationdate)
    /Trapped /False
}

%%%%%%%%%%%%%%%%%%%%%%%%%%%%%%%%%%%%%%%%%%%%%%%%%%%%%%%%%%%%%%%%%%%%%%%%%%%%%%%
%%%%%%%%%%%%%%%%%%%%%%%%%%%%%%%%%%%%%%%%%%%%%%%%%%%%%%%%%%%%%%%%%%%%%%%%%%%%%%%

\begin{document}
\selectlanguage{slovene}
\frontmatter
\setcounter{page}{1}
\renewcommand{\thepage}{}
\newcommand{\sn}[1]{"`#1"'} % slovenski narekovaji: \sn{besedilo}

%%%%%%%%%%%%%%%%%%%%%%%%%%%%%%%%%%%%%%%%
% naslovnica
 \thispagestyle{empty}%
  \begin{center}
    {\large\sc Univerza v Ljubljani\\%
      Fakulteta za računalništvo in informatiko\\%
     }
    \vskip 10em%
    {\autfont \tauthor\par}%
    {\titfont \ttitle \par}%
    {\vskip 3em \textsc{DIPLOMSKO DELO\\[5mm]
    VISOKOŠOLSKI STROKOVNI ŠTUDIJSKI PROGRAM\\ PRVE STOPNJE\\ RAČUNALNIŠTVO IN INFORMATIKA}\par}

    \vfill\null%
    {\large \textsc{Mentor}: \tmentor\par}

    {\vskip 2em \large Ljubljana, \tyear \par}
\end{center}

%%%%%%%%%%%%%%%%%%%%%%%%%%%%%%%%%%%%%%%%
% copyright stran
%%%%%%%%%%%%%%%%%%%%%%%%%%%%%%%%%%%%%%%%
\newpage
\thispagestyle{empty}

\vspace*{5cm}
{\small \noindent
  To delo je ponujeno pod licenco \textit{Creative Commons Priznanje avtorstva-Deljenje pod enakimi pogoji 2.5 Slovenija} (ali novejšo različico).
  To pomeni, da se tako besedilo, slike, grafi in druge sestavine dela kot tudi rezultati diplomskega dela lahko prosto distribuirajo,
  reproducirajo, uporabljajo, priobčujejo javnosti in predelujejo, pod pogojem, da se jasno in vidno navede avtorja in naslov tega
  dela in da se v primeru spremembe, preoblikovanja ali uporabe tega dela v svojem delu, lahko distribuira predelava le pod
  licenco, ki je enaka tej.
  Podrobnosti licence so dostopne na spletni strani \href{http://creativecommons.si}{creativecommons.si} ali na Inštitutu za
  intelektualno lastnino, Streliška 1, 1000 Ljubljana.

  \vspace*{1cm}
  \begin{center}
    {\Large\ccbysa}
  \end{center}
}

\vspace*{1cm}
{\small \noindent
  Izvorna koda diplomskega dela, njeni rezultati in v ta namen razvita programska oprema je ponujena pod licenco GNU General Public License,
  različica 3 (ali novejša). To pomeni, da se lahko prosto distribuira in/ali predeluje pod njenimi pogoji.
  Podrobnosti licence so dostopne na spletni strani \url{http://www.gnu.org/licenses/}.

  \medskip
  \noindent Repozitorij: \url{https://github.com/dilan23Jill/Analiza-razprav}
}

\vfill
\begin{center}
  \ \\ \vfill
  {\em
  Besedilo je oblikovano z urejevalnikom besedil \LaTeX.}
\end{center}

% prazna stran
\clearemptydoublepage

\thispagestyle{empty}
\
\vfill

\bigskip
\noindent\textbf{Kandidat:} \tauthor\\
\noindent\textbf{Naslov:} \ttitle\\
\noindent\textbf{Vrsta naloge:} Diplomska naloga na visokošolskem strokovnem programu prve stopnje Računalništvo in informatika \\
\noindent\textbf{Mentor:} \tmentor

\bigskip
\noindent\textbf{Opis:}\\
\hspace{1em}V diplomski nalogi izdelajte sistem, ki na podlagi vnesenega posnetka ustvari strukturiran pregled argumentacije. Sistem naj podpira nastop enega govorca in razpravo ena proti ena. Iz posnetka naj izdela prepis in pri tem jasno loči govorce. Veliki jezikovni modeli naj samodejno izluščijo argumente in jih zapišejo kot premise in sklep. Izrečene preverljive trditve naj sistem prek vnaprej določenih mehanizmov ovrednoti kot resnične, delno
resnične, zavajajoče, neresnične ali nepreverljive. Moderator ali drug stranski govorec naj bo zaznan in izločen iz ocenjevanja. Vsaka zaznana zmota naj ima dobesedni navedek, vsak izid preverjanja pa navedene vire. Sistem naj išče tudi logične zmote, razdeljene na formalne, neformalne in šibko sklepanje. Isti posnetek analizirajte večkrat in rezultate primerjajte.

\bigskip
\noindent\textbf{Title:} \ttitleEn

\bigskip
\noindent\textbf{Description:}\\
In your thesis, create a system that creates a structured overview of the argumentation based on the entered recording. The system should support the performance of a single speaker and a one-on-one discussion. It should produce a transcript from the recording, clearly distinguishing the speakers. Large language models should automatically extract arguments and record them as premises and conclusions. The system should evaluate the stated verifiable claims as true, partially true, misleading, false or unverifiable through predefined mechanisms. Depending on the results of the verification and the validity of the arguments, the model should return a categorical assessment, and the code should calculate the accuracy figures. The moderator or other external speaker should be detected and excluded from the evaluation. Each detected error should have a verbatim citation, and each verification result should have the sources listed. The system should also look for logical errors, divided into formal, informal and weak reasoning. Analyze the same recording several times and compare the results.

\vfill

\vspace{2cm}
\clearemptydoublepage

%%%%%%%%%%%%%%%%%%%%%%%%%%%%%%%%%%%%%%%%
% zahvala  (neobvezno — če je ne želiš, zakomentiraj do naslednje vrstice \clearemptydoublepage)
\thispagestyle{empty}\mbox{}\vfill\null\it%
\noindent
Iskreno se zahvaljujem mentorju izr.\ prof.\ dr.\ Slavku Žitniku za usmerjanje,
strokovno vodenje, ideje in hitro odzivnost.
Zahvaljujem se tudi družini, punci in prijateljem za vso spodbudo ter podporo v času
študija in pri nastajanju diplomskega dela.

\rm\normalfont
\clearemptydoublepage

%%%%%%%%%%%%%%%%%%%%%%%%%%%%%%%%%%%%%%%%
% posvetilo (neobvezno — celoten blok lahko izbrišeš)
%\thispagestyle{empty}\mbox{}{\vskip0.20\textheight}\mbox{}\hfill\begin{minipage}{0.55\textwidth}%
%% <-- posvetilo
%\normalfont\end{minipage}
%\clearemptydoublepage

%%%%%%%%%%%%%%%%%%%%%%%%%%%%%%%%%%%%%%%%
% kazalo
\pagestyle{empty}
\def\thepage{}
\tableofcontents{}

\clearemptydoublepage

\addcontentsline{toc}{chapter}{Povzetek}
\chapter*{Povzetek}

\noindent\textbf{Naslov:} \ttitle
\bigskip

\noindent\textbf{Avtor:} \tauthor
\bigskip

\noindent
Diplomska naloga obravnava izdelavo sistema, ki
posnetek razprave samodejno pretvori v strukturiran pregled predstavljene argumentacije. Spletne razprave pogosto trajajo več ur, med tem pa je veliko rečenega nepomembnega, ali pa gre za osebne napade. Gledalec, ki bi hotel presoditi, kdo je katero trditev dejansko podprl in kdo se je izmikal, mora za to porabiti veliko časa. Ročna analiza argumentacije je zamudna, sodba je pa lahko pristranska. Izdelali smo sistem, ki celotno pot od videoposnetka do berljivega poročila opravi samodejno. Posnetek prenese, naredi prepis z ločevanjem govorcev in argumentacijo analizira v zaporednih korakih, pri razpravi dveh govorcev v petih, pri nastopu enega govorca v štirih. Iz prepisa najprej izlušči argumente s premisami, nato v njih poišče logične zmote, preverljive podatke iz premis preveri v znanstvenih in spletnih virih, poveže odgovore z napadenimi argumenti in napiše povzetek. Osrednja zasnovna odločitev je ločnica med presojo in izračunom. Jezikovni model vrne izključno nominalne ocene iz vnaprej določenih množic in ne poda nobene številske ocene. Številke v poročilu nastanejo v sistemu, ki nominalne ocene prešteje, poveže med koraki in prikaže.

\bigskip

\noindent\textbf{Ključne besede:} \tkeywords.
\clearemptydoublepage

%%%%%%%%%%%%%%%%%%%%%%%%%%%%%%%%%%%%%%%%
% abstract
\selectlanguage{english}
\addcontentsline{toc}{chapter}{Abstract}
\chapter*{Abstract}

\noindent\textbf{Title:} \ttitleEn
\bigskip

\noindent\textbf{Author:} \tauthor
\bigskip

\noindent
The thesis deals with the creation of a system that automatically converts a recording of a discussion into a structured overview of the presented argumentation. Online discussions often last for several hours, during which much of what is said is irrelevant or involves personal attacks. A viewer who wants to judge who actually supported which claim and who evaded it must spend a lot of time doing so. Manual analysis of the argumentation is time-consuming, and the judgment can be biased. We built a system that automatically completes the entire journey from video to readable report. It downloads the recording, produces a transcript with speaker separation and analyzes the argumentation in successive steps, five for a debate between two speakers and four for a single speaker. It first extracts arguments with their premises, then names the logical fallacies in them, checks the verifiable facts in the premises against scientific and online sources, links rebuttals to the arguments they attack and writes a summary. The central design decision is the dividing line between judgment and calculation. The language model returns exclusively categorical ratings from predefined sets and does not provide any numerical ratings. The numbers in the report are computed by the system, which counts the categorical ratings, links them across steps and displays them.

\bigskip

\noindent\textbf{Keywords:} \tkeywordsEn.
\selectlanguage{slovene}
\clearemptydoublepage

%%%%%%%%%%%%%%%%%%%%%%%%%%%%%%%%%%%%%%%%
\mainmatter
\setcounter{page}{1}
\pagestyle{fancy}

%  1  UVOD

\chapter{Uvod}
\label{ch:uvod}

Javne razprave o pomembnih vprašanjih so danes večinoma dostopne kot dolgi
videoposnetki. Kdor jih spremlja, hitro pozabi, kaj natanko je kdo trdil in s
katerimi razlogi je to podprl. Poslušanje vzame več ur, sproti pa je težko
preveriti, ali navedeni podatki držijo in ali sklep res sledi iz razlogov, ki so
bili zanj navedeni. Kdor si hoče ustvariti mnenje, mora posnetek pogledati do
konca in si vzporedno delati zapiske.

Od tod izhaja zamisel te naloge. Potreben je pripomoček, ki posnetek razprave
prebere namesto gledalca in pokaže, kaj je kdo zagovarjal, s čim je to podprl,
kje se je pri sklepanju zmotil in katere izrečene podatke potrjujejo zunanji
viri. Bralec namesto večurnega poslušanja dobi pregled, ki ga prebere v nekaj
minutah, skupaj s povezavami do virov, če želi kaj preveriti sam. Takšno orodje
je uporabno za vsakogar, ki nima časa za celoten posnetek ali pa raje bere, kot
posluša.

Cilj naloge je izdelati sistem, ki celotno pot opravi samodejno. Sistem sprejme
povezavo do posnetka ali naloženo datoteko, izdela prepis in ga razdeli po
govorcih. Podpira razpravo dveh govorcev in nastop enega govorca. Iz prepisa nato z velikimi jezikovnimi modeli izlušči argumente, torej sklepe govorcev in premise, na katerih stojijo. Vsak razlog, ki
navaja preverljiv podatek, sooči z zunanjimi viri, v sklepanju pa poišče
morebitne logične zmote. Izid je pregleden zapis argumentacije vseh sodelujočih,
dostopen v spletni aplikaciji in izvozljiv v obliki PDF. Videz aplikacije prikazujejo naslednje slike.

\begin{figure}[H]
\centering
\includegraphics[width=1\textwidth,keepaspectratio]{nova_analiza.png}
\caption{Stran Nova analiza, na kateri uporabnik vnese povezavo ali datoteko in
izbere možnosti analize.}
\label{fig:nova-analiza}
\end{figure}

\begin{figure}[H]
\centering
\includegraphics[width=\textwidth,height=0.60\textheight,keepaspectratio]{pregled.png}
\caption{Končana analiza v spletnem vmesniku. Na vrhu je glava z nastavitvami in navigacijo, pod njo možnosti za urejanje, ponovno analizo, izvoz v PDF in pojasnilo, kako brati analizo, spodaj pa časovnica argumentov in pregled dejstev.}
\label{fig:izdelek}
\end{figure}


%  2  OZADJE IN SORODNA DELA

\chapter{Ozadje in sorodna dela}
\label{ch:ozadje}

Naloga zajema področja teorije argumentacije, zaznavo
logičnih zmot, samodejno rudarjenje argumentov, samodejno preverjanje dejstev in uporabo velikih jezikovnih modelov v vlogi ocenjevalcev. To poglavje predstavi pojme in dosedanje delo na teh področjih, na koncu pa utemelji, v čem se opisani pristop razlikuje.

\section{Osnovni pojmi}
\label{sec:pojmi}

\textbf{Trditev} je stavek, ki nekaj zatrjuje. Lahko je objektiven ali subjektiven. Primer: \sn{Zunaj je lepo vreme.} Definicijo povzemamo po klasični logiki, kjer je trditev nosilec resničnostne vrednosti~\cite{copi2014}. Potrebujemo jo, ker so trditve osnovna enota, iz katere sistem gradi argumente. 

\textbf{Dejstvena trditev} je objektivna trditev, ki jo je mogoče soočiti z zunanjimi viri, ker govori o svetu: navaja število, datum, raziskavo, merljivo vrednost ali to, kdo je kaj rekel oziroma storil. Trditvi, ki pove, kaj bi bilo treba storiti ali kaj je pravično, ne pravimo dejstvena, saj izraža subjektivno mnenje, ki ga z viri ni mogoče potrditi ali ovreči. Primer dejstvene trditve je \sn{stopnja brezposelnosti se je lani znižala za dve odstotni točki}, primer subjektivne pa \sn{takšna politika je nepravična}. Razlikovanje povzemamo po literaturi o samodejnem preverjanju dejstev, kjer je prvi korak izbira trditev, ki jih je mogoče preveriti~\cite{guo2022survey,hassan2017claimbuster}. Za sistem je pomembno, ker preverjanje dejstev obravnava samo dejstvene trditve.

\textbf{Premisa} je osnovni gradnik argumenta, ki vsebuje neko trditev in iz katere sledi sklep~\cite{copi2014}. V sistemu premisa poleg trditve vsebuje še razlog, ki ga je zanjo navedel govorec (razdelek~\ref{ssec:prehod1}).

\textbf{Sklep} je izpeljava iz premis. Primer: iz premis \sn{Vsi ljudje umrejo} in \sn{Sokrat je človek} sledi sklep \sn{Sokrat bo umrl}. Definicijo povzemamo po klasični logiki, kjer je sklep tisto, kar je v argumentu potrjeno na podlagi premis~\cite{copi2014}. 

\textbf{Argument} je sestava premis in sklepa. 

\textbf{Zmota} je napaka v sklepanju, zaradi katere premise sklepa ne podpirajo ali ga podpirajo šibkeje, kot trdi govorec. Primer: sklep, da Spartanci niso bili najboljši bojevniki, ker bi jih premagal rimski legionar, primerja vojski iz različnih obdobij, kar je lažna enakovrednost. 

\textbf{Razsodba} je izid preverjanja ene dejstvene trditve. Razsodbe so \sn{resnično}, \sn{delno resnično}, \sn{zavajajoče}, \sn{neresnično} in \sn{nepreverljivo}. Primer: trditev, da je Triglav visok 2864 metrov, dobi razsodbo \sn{resnično}. Izraz povzemamo po literaturi o preverjanju dejstev, kjer je razsodba izid zadnjega koraka~\cite{guo2022survey}. 

\textbf{Zavrnitev} je odgovor enega govorca, s katerim spodbija argument drugega, bodisi tako, da pokaže na zmoto v sklepanju, bodisi tako, da popravi dejstva. Primer: govorec trdi, da so bili Spartanci najboljši bojevniki antike, ker je pri Termopilah 300 Spartancev zadrževalo perzijsko vojsko. Drugi govorec lahko to premiso zavrne na dva načina. Lahko pokaže na šibko sklepanje, saj iz ene same bitke ni mogoče sklepati o vojaški moči v celotnem obdobju, kar je prehitra posplošitev. Lahko pa popravi dejstvo, saj pri Termopilah ni bilo le 300 Spartancev, ampak nekaj tisoč grških vojakov. Zavrnitev ustreza razmerju napada med argumentacijskimi enotami v rudarjenju argumentov~\cite{stede2018}.

\textbf{Izmikanje} je primer, ko govorec na neposredno vprašanje nasprotnika ne odgovori ali odgovori le delno. Primer: na vprašanje, ali bi Ahil premagal Herkula, govorec odgovori le, da je Herkul bog, in preide na drugo temo. Opredelitev je lastna in prilagojena nalogi.

\textbf{Razprava} je izmenjava nasprotujočih si mnenj, v kateri sodelujeta vsaj dva udeleženca. Vsak svoje stališče utemeljuje z razlogi in poskuša izpodbiti argumente nasprotnika. Pri formalnih razpravah sodeluje še moderator, ki gledalcem predstavi potek in skrbi, da udeleženca ostaneta pri temi. Opredelitev se naslanja na Waltonov opis prepričevalnega dialoga~\cite{walton2008}. 

\textbf{Stališče} je glavna teza, ki jo govorec zagovarja v celotni razpravi. Argumenti govorca stališče večinoma podpirajo, ni pa nujno. Govorec ima eno stališče in lahko več argumentov. Primer: stališče \sn{Trojanska vojna, kot jo opisuje Homer, se ni zgodila.} podpira argument \sn{Homerjeva vojna ni zgodovinski dogodek, saj izkopavanja kažejo uničenje, ne pa desetletnega obleganja.} Razlikovanje potrebujemo, ker sistem za vsakega govorca zapiše stališče in seznam argumentov zanj.

\textbf{Poziv} je besedilo, ki ga jezikovni model prejme kot vhod ob klicu. Sestavljen je iz sistemskega poziva, ki določi vlogo modela in pravila, ter uporabniškega poziva, ki vsebuje nalogo in gradivo za ta klic. Primer: pri zaznavi zmot sistemski poziv modelu pove, naj poroča nevtralno, ne razglasi zmagovalca in zmoto zapiše le, če pokaže premiso z napako, uporabniški poziv pa vsebuje nalogo \sn{Za vsak spodnji argument presodi, ali sklep sledi iz premis in ali vsebuje poimenovano zmoto} in za njo izluščene argumente. 

\section{Zgradba argumenta}
\label{sec:zgradba-argumenta}

Sistem gradi na modelu angleškega filozofa Stephena Toulmina~\cite{toulmin1958}, ker je primeren za razčlenjevanje argumentov iz vsakdanjih razprav. Model argument razčleni na šest delov: sklep (angl.\ claim), premise, ki ga podpirajo (angl.\ data), sklepanje, ki premise povezuje s sklepom (angl.\ warrant), oporo za samo sklepanje (angl.\ backing), omejitev veljavnosti (angl.\ qualifier) in zavrnitev (angl.\ rebuttal).

Toulmin z izrazom \textit{rebuttal} sicer označi okoliščine, v katerih sklep ne velja in jih navede govorec sam. Prilagoditve modela v rudarjenju argumentov pa z njim označijo nasprotno stališče, ki sklep napade~\cite{habernal2017}. V tej nalogi ga uporabljamo v drugem pomenu, torej kot odgovor nasprotnika.

Kako so ti deli zajeti v sistemu, prikazuje tabela~\ref{tbl:toulmin}.

\begin{table}[H]
\centering
\small
\arrayrulecolor{crta}
\begin{tabular}{@{}p{0.24\textwidth}p{0.13\textwidth}p{0.49\textwidth}@{}}
\toprule
\textbf{Del} & \textbf{Lastno polje} & \textbf{Kje v sistemu nastopa} \\
\midrule
sklep (\textit{claim}) & da & polje \texttt{argument}, sklep z jedrom utemeljitve \\
\midrule
premise (\textit{data}) & da & seznam premis argumenta \\
\midrule
sklepanje (\textit{warrant}) & ne & zmota šibkega sklepanja, vrsta zavrnitve izpodbijanje sklepanja \\
\midrule
opora (\textit{backing}) & ne & sistem je ne zajame \\
\midrule
omejitev (\textit{qualifier}) & ne & zmota prehitre posplošitve \\
\midrule
zavrnitev (\textit{rebuttal}) & da & odgovor nasprotnika \\
\bottomrule
\end{tabular}
\caption{Preslikava Toulminovega modela v podatkovni model sistema.}
\label{tbl:toulmin}
\end{table}

\section{Logične zmote}
\label{sec:zmote}
Logična zmota je napaka v sklepanju, zaradi katere argument ne dokazuje tega, kar naj bi dokazoval, čeprav se lahko zdi prepričljiv~\cite{copi2014}. Formalne zmote so napake v obliki sklepanja in jih je mogoče prepoznati brez poznavanja vsebine. Primer je potrjevanje posledice: \sn{Če dežuje, so ceste mokre. Ceste so mokre, torej dežuje.} Sklep ne sledi, ker so ceste lahko mokre tudi zaradi drugega vzroka. Neformalne zmote, ki so v javnih razpravah bistveno pogostejše, so odvisne od konteksta in jih brez poznavanja tematike ni lahko prepoznati~\cite{walton2008}. Primer je napad na osebo: \sn{Zdravnik pravi, da je kajenje škodljivo, a sam kadi, zato mu ne moremo verjeti.} Namesto trditve je napaden govorec. Tretja skupina je šibko sklepanje, ki se pojavi, ko premise sklep podpirajo šibkeje, kot trdi govorec, ali iz njih ne sledi. Primer je prehitra posplošitev: \sn{Dvakrat sem jedel v tej restavraciji in obakrat je bila hrana slaba, torej je tam hrana vedno slaba.} Dve izkušnji sklep podpirata, a ne v tolikšni meri.

Zaznava zmot v govorjenih razpravah je obdelana tudi v literaturi. Mancini in sodelavci so zbrali posnetke ameriških predsedniških soočenj in v njih ročno označili odseke, ki vsebujejo zmoto, ter za vsak odsek zapisali, katera od šestih zmot je v njem. Taka zbirka označenih primerov se imenuje korpus in služi za učenje ter preizkušanje modelov. Razvrščanje je večmodalno, kar pomeni, da model poleg prepisa dobi tudi zvok odseka. Ugotovili so, da zvok izboljša mero $F_1$ za štiri do osem odstotnih točk~\cite{mancini2024fallacy}. Mera $F_1$ je uveljavljena mera uspešnosti razvrščanja in združuje dvoje: natančnost, torej kolikšen delež označenih primerov je bil res pravilen, in priklic, torej kolikšen delež pravilnih primerov je sistem sploh našel. Zavzame vrednosti med $0$ in $1$, pri čemer višja pomeni boljši izid. Njihov sistem pa dobi odsek, za katerega je že znano, da vsebuje zmoto, in izbere le, katera zmota je v njem. Ne pove torej, kje v razpravi je zmota.

\section{Samodejno rudarjenje argumentov}
\label{sec:argument-mining}

Rudarjenje argumentov je samodejno prepoznavanje in izluščanje zgradbe sklepanja, izraženega v naravnem jeziku~\cite{lawrence2019argument}. Naloga se običajno razčleni na tri korake: razdelitev besedila na argumentacijske enote, razvrstitev enot na sklepe in premise, ter napoved razmerij med njimi, torej podpore ali napada~\cite{stede2018}.

Govorjena razprava je za razčlenjevanje argumentov najtežji vhod. Povedi ostajajo nedokončane, govorci si segajo v besedo, meja med samostojnim argumentom in podtočko istega argumenta pa je pogosto stvar presoje. 

Kako težka je naloga, pokaže ujemanje med ljudmi, ki isto besedilo označujejo neodvisno drug od drugega. Mere ujemanja, kot so $\pi$, $\kappa$ in $\alpha$, povedo, koliko se označevalci strinjajo bolj, kot bi se ujeli že po naključju. Vrednost $1$ pomeni popolno ujemanje, $0$ pa ujemanje na ravni naključja. V jezikoslovnih nalogah vrednosti nad $0{,}80$ veljajo za zanesljive, vrednosti pod $0{,}60$ pa za šibke. Pri prepričevalnih esejih, torej šolskih spisih, v katerih pisec zagovarja stališče, so trije označevalci dosegli $\pi = 0{,}66$ za sklepe in $\pi = 0{,}71$ za premise~\cite{stab2014}. V spletnih razpravnih nitih, kjer je besedilo manj urejeno, je ujemanje padlo na $\kappa = 0{,}53$ oziroma $\kappa = 0{,}55$~\cite{collagree2018}.

Z velikimi jezikovnimi modeli se je pristop k temu področju spremenil, saj lahko en sam poziv opravi več naštetih podnalog hkrati. Pregled Lija in sodelavcev med odprtimi izzivi navaja sklepanje čez dolg kontekst, razložljivost izida in pomanjkanje označenih podatkov~\cite{li2025llmargmining}.

Govorjenim razpravam je namenjen korpus M-Arg~\cite{mestre2021marg}. Zajema pet soočenj ameriških predsedniških volitev 2020, sedem ur zvoka s prepisi in 6527 izjav. Spletni delavci, ki so za označevanje najeti prek spleta, so za 4104 pare izjav določili, ali prva izjava drugo podpira, napada ali z njo ni v nobenem razmerju. Takratni modeli dosežejo skupno točnost 0,86, vendar je zbirka močno neuravnotežena, saj večina parov ni v nobenem razmerju. Model, ki skoraj vedno odgovori, da razmerja ni, zato doseže visoko točnost, čeprav napadov in podpor ne najde. Mera $F_1$ za napade znaša 0,24 in za podpore 0,21, visoka točnost pa izvira predvsem iz pravilnega prepoznavanja parov brez razmerja. Ujemanje med označevalci je bilo pri tem nizko, $\alpha = 0{,}43$. Korpus poleg tega razmerja med izjavami le razvršča, zgradbe argumenta s premisami in sklepom pa ne izlušči. 

\section{Samodejno preverjanje dejstev}
\label{sec:fact-checking-sorodna}
Samodejno preverjanje dejstev ugotavlja resničnost trditve s pomočjo obdelave naravnega jezika in zunanjih virov znanja~\cite{guo2022survey}. Razdeli se na
štiri korake, in sicer zaznavo preverljivih trditev, pridobivanje virov, tvorjenje razsodbe in tvorjenje utemeljitve.

Zaznava trditev mora ločiti preverljivo dejstveno trditev od zagovarjanja stališča. Pridobivanje virov je v obstoječih merilnih zbirkah pogosto omejeno na en sam vir, najpogosteje Wikipedijo, kar poenostavi vrednotenje, a slabo služi v resničnih razpravah, kjer se trditve raztezajo od medicine do aktualne politike. Tvorjenje razsodbe se zaplete pri sestavljenih trditvah, pri katerih
je en del točen, drugi pa ne.

Najbolj uveljavljen sklenjen sistem je ClaimBuster~\cite{hassan2017claimbuster}, prvi, ki je celotno pot od besedila do razsodbe izvedel samodejno. Vsakemu stavku dodeli oceno, kako verjetno vsebuje trditev, vredno preverjanja. Stavek \sn{Brezposelnost je lani padla na pet odstotkov} tako dobi visoko oceno, stavek \sn{Hvala vsem, da ste prišli} pa nizko. Pri tem doseže natančnost 79 in priklic 74 odstotkov. Za izbrane stavke nato poišče že objavljena preverjanja in za nove trditve tvori poizvedbe po spletu. Preizkušen je bil v živo med soočenji ameriških volitev leta 2016. Deluje pa na ravni posameznega stavka, saj ne pove, čigav argument preverjena trditev podpira, in ne presodi, ali sklep iz nje sledi.

\section{Veliki jezikovni modeli kot ocenjevalci}
\label{sec:llm-judge}

Veliki jezikovni modeli se vse pogosteje uporabljajo namesto človeških ocenjevalcev. Pri tem so znane težave: model daje prednost daljšim odgovorom in bolje ocenjuje besedila, ki jih je sam ustvaril~\cite{zheng2023judging}. Ko izbira med dvema odgovoroma, lahko izbere drugega samo zato, ker sta v pozivu zamenjala mesti. Raziskava petnajstih modelov v vlogi ocenjevalca na več kot 150.000 primerih je pokazala, da pristranskost glede na položaj odgovora v pozivu ni naključna, ampak se sistematično razlikuje med modeli in nalogami~\cite{shi2025judging}. Zanesljivost ocenjevalcev je zato
samostojna raziskovalna tema~\cite{gu2024}. Ker presoja modela ni zanesljiva, smo jo omejili na izbiro iz zaprte množice, v kateri je vsaka vrednost opredeljena s pravilom. 

Pet razsodb smo izbrali na podlagi lastnega znanja in rednega spremljanja širokega nabora razprav v zadnjih letih. Izbrane so tako, da vsaka pokrije drugačen izid preverjanja. \sn{Resnično} pomeni, da je trditev točna, kot je bila izrečena, \sn{neresnično} pa, da je v nasprotju z zanesljivimi viri. Ti dve vrednosti ne zadoščata, ker veliko trditev v razpravah ni ne povsem točnih ne povsem napačnih. \sn{Delno resnično} zato pokrije trditev, katere smer je pravilna, a se podrobnosti, na primer število, razlikujejo. \sn{Zavajajoče} pokrije trditev, katere posamezni deli so resnični, a kot celota ustvarja napačen vtis. Od delno resnične je ločena, ker težava ni v netočnosti, ampak v vtisu, ki ga trditev pusti pri poslušalcu. \sn{Nepreverljivo} pomeni, da ustreznih virov ni, in jo poziv predpiše za primer, ko gradivo ne da podlage za razsodbo.

\section{Umestitev naloge glede na sorodna dela}
\label{sec:umestitev}

Obstoječi pristopi rešujejo posamezne dele naloge, ki si jo zastavlja to delo, nobeden pa ne pokriva poti od posnetka do poročila v celoti. M-Arg zajema govorjene razprave, a razmerja med izjavami le razvršča. Mancini in sodelavci zaznavajo zmote nad že označenimi odseki. ClaimBuster preverja dejstva, a na ravni posameznega stavka in brez zgradbe argumenta. Project Debater~\cite{slonim2021} argumente tvori in v razpravi sodeluje, pri opisanem sistemu pa je naloga obrnjena, saj argumente prepoznava in jih ne tvori.

Sorodna dela se razlikujejo tudi po vhodu. Nobeno od naštetih del ne prejme samo zvoka. ClaimBuster prejema podnapise televizijskega prenosa, M-Arg in korpus zmot pa se opirata na že pripravljene prepise z oznakami govorcev. Ločevanje govorcev je pri njih dano, v sistemu pa je del cevovoda.

Sistem iz povezave naloži posnetek, iz njega naredi prepis z ločevanjem govorcev, izlušči argumente s premisami, najde in poimenuje zmote, preveri trditve iz premis v več virih in poveže zavrnitve z napadenimi argumenti. Celotno pot prikazujejo slike~\ref{fig:cevovod}--\ref{fig:cevovod3}.

Prispevek tega dela je torej cenovno dostopen sistem s celotno potjo obdelave od posnetka do poročila v eni aplikaciji, v kateri je vsaka navedba povezana z argumentom, iz katerega je nastala, in z viri, ki jih uporabnik lahko sam odpre.

\chapter{Zasnova sistema}
\label{ch:zasnova}

V tem poglavju so opisane zahteve in ključne odločitve, ki so oblikovale sistem. Obdelava posnetka je razdeljena na več delov, pri izbiri modela za prepis in modelov za posamezne korake analize pa smo upoštevali ceno in točnost. Iskali smo točko, kjer za najnižjo ceno dobimo zadostno dober rezultat. To smo dosegli z testiranjem in primerjanjem vsakega koraka analize.  Razdelek~\ref{sec:zahteve} navede funkcionalne in nefunkcionalne zahteve, razdelek~\ref{sec:arhitektura} pa opiše razdelitev aplikacije na odjemalca, strežnik in opravila v ozadju. Razdelek~\ref{sec:cevovod} poda pregled cevovoda. Sledi opis odločitev o tem, da model vrne le nominalne vrednosti, vse številke pa izračuna sistem (razdelek~\ref{sec:kategorije-vs-stevilke}), in da je analiza razdeljena na več zaporednih korakov (razdelek~\ref{sec:koraki}). To poglavje utemelji, zakaj je sistem zasnovan tako, kako so posamezni koraki izvedeni, pa opiše poglavje~\ref{ch:implementacija}.

\section{Zahteve}
\label{sec:zahteve}

Funkcionalne zahteve so:

\begin{itemize}
\item sprejem povezave do videoposnetka na YouTubu ali neposredno naložene
datoteke, z možnostjo izbire časovnega odseka,
\item prepis z ločenimi govorci, da je razvidno, kaj je povedal posamezni
govorec,
\item preverjanje dejstvenih trditev z razsodbami in navedenimi viri,
\item analiza argumentacije, torej argumenti s premisami, preslikava odgovorov (zavrnitev
in izmikanj), zaznane logične zmote ter končna sinteza,
\item spletna aplikacija, ki omogoča ročno popravljanje analize in ponovitev
analize istega posnetka brez ponovnega prepisa,
\item izvoz poročila v obliki PDF.
\end{itemize}

Nefunkcionalne zahteve so:

\begin{itemize}
\item Sledljivost: vsaka zaznana zmota mora navesti argument in njegov
del, v katerem je napaka, vsaka razsodba preverjanja dejstev pa vire in premiso,
ki jo preverja (razdelek~\ref{sec:podatkovni-model}).
\item Ponovljivost: ponovni zagon analize na istem posnetku mora vrniti
vsebinsko primerljiv rezultat, kar pokaže, da je sistem konsistenten. 
\item Nevtralnost: sistem ne sme kaznovati stališča, ampak kvečjemu
način njegove utemeljitve. Poziv za zaznavo zmot zato od modela zahteva, da navede premiso, ki nosi napako, in poimenuje napako v zgradbi sklepanja, pri dvoumnem primeru pa mora obrazložitev navesti obe branji, bralec pa lahko zaznavo
potrdi, zavrne ali odstrani.
\item Obvladljiv strošek: analiza sprejetih posnetkov mora ostati v
meji dveh evrov, sicer sistem ni uporaben za redno rabo.
\item Odpornost: jezikovni modeli občasno vrnejo okvarjen ali nepopoln
izhod, zunanje storitve pa omejujejo število klicev, zato mora cevovod take
napake preživeti brez izgube celotne analize (razdelek~\ref{sec:determinizem}).
\end{itemize}

\section{Arhitektura aplikacije}
\label{sec:arhitektura}

\begin{figure}[H]
\centering
\includegraphics[width=\textwidth]{arhitektura.pdf}
\caption{Arhitektura aplikacije. Odjemalec pošlje asinhrono zahtevo za analizo, strežnik takoj vrne oznako opravila in zažene opravilo v ločeni niti. Opravilo kliče zunanje storitve, izid in prepis pa shrani v hrambo. Odjemalec medtem s sinhronimi zahtevami poizveduje po stanju opravila.}
\label{fig:arhitektura}
\end{figure}

Aplikacija je razdeljena na odjemalca, strežnik in zunanje storitve (slika~\ref{fig:arhitektura}). Odjemalec je enostranska aplikacija v ogrodju React, ki teče v brskalniku. Strežnik je napisan v jeziku Python z ogrodjem FastAPI in odjemalcu ponuja vmesnik REST.

Odjemalec pošlje asinhrono zahtevo za analizo, strežnik ustvari opravilo in
vrne njegovo oznako, cevovod pa steče v ločeni niti. Zahteva je asinhrona, ker odgovor ne vsebuje izida analize, ampak le oznako opravila, s katero ga odjemalec pridobi pozneje. Pred tem strežnik preveri prijavo uporabnika in dolžino posnetka, ki ne sme presegati 23 minut.

Nit med izvajanjem sproti zapisuje stanje opravila in korak, ki trenutno teče. Odjemalec medtem v kratkih presledkih pošilja sinhrono zahtevo za stanje, ki takoj vrne trenutno stanje, in ga prikazuje uporabniku (slika~\ref{fig:loading}). Ko je opravilo končano, odjemalec prikaže rezultat.

Izid končane analize se shrani v podatkovno bazo SQLite, v kateri so tudi uporabniki in njihove seje. Prepisi končanih analiz se hranijo ločeno od rezultatov, kot datoteke .txt na strežniku. Ponovna analiza istega posnetka zato preskoči prenos in prepis (razdelek~\ref{ssec:rerun}).

Cevovod v niti kliče zunanje storitve. Zvok posnetka prenese orodje yt-dlp, prepis z ločevanjem govorcev izdela model OpenAI, korake analize opravijo jezikovni modeli ponudnikov Anthropic in OpenAI, preverjanje dejstev pa gradivo zbira iz iskalnikov znanstvenih in spletnih virov.

Odjemalec strukturirane rezultate prikaže po govorcih in argumentih, omogoča ročno popravljanje analize ter izvoz poročila v PDF
(razdelek~\ref{sec:aplikacija}). Končne točke vmesnika REST navaja razdelek~\ref{ssec:backend}.

\section{Pregled sistema}
\label{sec:cevovod}

Sistem je zasnovan kot cevovod zaporednih korakov, in sicer prenos zvoka, prepis z ločevanjem govorcev, večkoračna analiza argumentacije, obdelava izhodov v kodi ter shranjevanje in prikaz. Celoten cevovod prikazujejo slike~\ref{fig:cevovod}--\ref{fig:cevovod3}. Za prepis smo izbrali model \texttt{gpt-4o-transcribe-diarize}, ki v enem klicu vrne besedilo in oznake govorcev, pri naših preizkusih pa je bil natančnejši od drugih modelov, tudi pri slovenskih posnetkih.

Ker je cevovod razdeljen na natančno določene zaporedne korake, je vsakega mogoče razvijati in preizkušati posebej, brez posega v preostanek aplikacije.

Ena od zasnovnih odločitev je, da preverjanje dejstev za vsako preverljivo
trditev zbere gradivo iz več vrst virov hkrati, kar poveča možnost, da se najde vir, ki trditev potrdi ali ovrže. Na koncu model nad vsem zbranim gradivom vrne eno razsodbo. Ta bralcu pove, ali je trditev najverjetneje resnična, ne da bi moral sam pregledati vsak vir. Vsak vir pa dobi še svojo oznako, kam kaže (razdelek~\ref{ssec:razsojanje}).

\subsection{Podatkovni model}
\label{sec:podatkovni-model}

Koraki analize si podajajo zapise v obliki JSON, katerih obliko določajo sheme (razdelek~\ref{ssec:validacija}). Izhod prvega koraka je slovar govorcev. Za vsakega govorca vsebuje stališče (\texttt{position}) in seznam argumentov (\texttt{arguments}), vsak argument pa sklep z jedrom utemeljitve (\texttt{argument}) in seznam premis (\texttt{premises}). Argument je osrednji zapis, na katerega se sklicujejo vsi poznejši koraki. Po prvem koraku mu koda dodeli oznako \texttt{arg\_id} oblike \texttt{govorec\#i}, kjer je \texttt{i} zaporedna številka argumenta pri tem govorcu (razdelek~\ref{ssec:arg-id}).

Zapisi poznejših korakov se nanj sklicujejo z oznako:
\begin{itemize}
\item zmota (drugi korak) nosi \texttt{arg\_id} argumenta, v katerem je najdena, in po potrebi \texttt{premise\_index}, številko premise z napako,
\item preverljiva trditev (tretji korak) nosi \texttt{arg\_id} in \texttt{premise\_index} premise, iz katere izvira,
\item zavrnitev (četrti korak) nosi \texttt{target\_arg\_id}, oznako napadenega argumenta nasprotnika.
\end{itemize}
Kako sistem ta sklicevanja preveri, opisuje razdelek~\ref{ssec:povezovanje}. Kako so ti zapisi povezani v vmesniku, prikazuje slika~\ref{fig:vmesnik-arg1}.

\begin{figure}[H]
\centering
\includegraphics[width=\textwidth,height=0.75\textheight,keepaspectratio]{arg1.png}
\caption{V odprtem argumentu zgoraj so oštevilčene premise in iz njih izpeljan sklep, ob vsaki premisi pa pika v barvi razsodbe za trditve, ki iz nje izvirajo. Pod argumentom so zaznane zmote s kategorijo in delom argumenta, v katerem je napaka. S klikom na gumba \sn{Correct} in \sn{Not correct} uporabnik potrdi ali zavrne zaznavo.}
\label{fig:vmesnik-arg1}
\end{figure}

\section{Nominalne vrednosti modela in izračun v kodi}
\label{sec:kategorije-vs-stevilke}

Model nikjer ne vrne gole ocene na lestvici, ampak na vsakem mestu, kjer poziv zahteva presojo, vrne vrednost z vnaprej pripravljenega seznama. Ker presoja modela ni zanesljiva (razdelek~\ref{sec:llm-judge}), smo jo omejili na izbiro iz zaprte množice, v kateri je vsaka vrednost opredeljena s pravilom, saj je taka izbira med zagoni bolj konsistentna. Ta mesta našteva tabela~\ref{tbl:zaprte}.

\begin{table}[H]
\centering
\small

\arrayrulecolor{crta}
\begin{tabular}{@{}p{0.30\textwidth}p{0.62\textwidth}@{}}
\toprule
Vrsta vrednosti & Dovoljene vrednosti \\
\midrule
skupna razsodba o trditvi & resnično, delno resnično, zavajajoče, neresnično, nepreverljivo \\
\midrule
kam kaže posamezen vir & resnično, delno resnično, zavajajoče, neresnično, nepreverljivo \\
\midrule
ime zmote & ena od dvaintridesetih v zaprtem slovarju ali \texttt{other} \\
\midrule
vrsta zavrnitve & neposredno nasprotovanje, rušenje premise, druga razlaga, izpodbijanje sklepanja \\
\midrule
vrsta izmikanja & preusmeritev, menjava teme, odgovor brez vsebine, delni odgovor, preglasitev \\
\midrule
tip trditve & statistična, zgodovinska, znanstvena, navedek, politična, zdravstvena, ekonomska, zemljepisna \\
\midrule
vloga udeleženca & glavni govorec, debater, moderator \\
\bottomrule
\end{tabular}
\caption{Mesta, kjer model vrne vrednost iz zaprte množice.}
\label{tbl:zaprte}
\end{table}

Nobena od teh vrednosti ni ocena kakovosti argumenta. Razsodbi povesta, ali izrečeni podatek drži, ena o trditvi kot celoti in druga o vsakem viru posebej. Zakaj je razsodb pet in kaj pomeni vsaka, pojasnjuje razdelek~\ref{sec:llm-judge}.
Ime zmote poimenuje napako v sklepanju. Vrsti zavrnitve in izmikanja opišeta,
kako je potekala izmenjava med govorcema. Tip trditve in vloga udeleženca zapise le razvrstijo, pri čemer tip trditve določi še, kateri
zbiralci virov stečejo (razdelek~\ref{ssec:engine-routing}).

Številke v poročilu izračuna sistem, ki te nominalne vrednosti prešteje. Takih
številk je pet:

\begin{itemize}
\item koliko trditev je bilo preverjenih,
\item koliko trditev je dobilo katero od petih razsodb, skupno in po govorcih,
\item koliko virov stoji za posamezno trditvijo,
\item koliko različnih domen ti viri pokrivajo,
\item koliko teh virov kaže v katero smer, torej seštevek oznak pri eni sami
trditvi.
\end{itemize}

Presoja modela je omejena na izbiro z zaprtega seznama, vse nadaljnje računanje
pa teče znotraj sistema in je zato ob istem vhodu vedno enako. Zamenjava modela s tem
spremeni kakovost presoje, ne pa načina obdelave, zato je učinek novega modela
mogoče izmeriti brez sprememb v kodi.

Sistem ne presoja, ali je argument dober, ali je po zavrnitvi obstal in kako je
pomemben. To ni njegova naloga, saj je sodba o tem, kdo je bil prepričljivejši,
odvisna od bralca. Sistem pokaže, kaj je bilo rečeno, in
poimenuje zmote v sklepanju, odločitev pa prepusti bralcu. Sinteza povzame, kaj je kdo zagovarjal, na čem je to utemeljil, kje se je
izmikal in katera vprašanja so ostala odprta, ne da bi izrekla, kdo je bil
boljši. Po korakih razčlenjeno ločnico med presojo modela in izračunom v kodi
navaja dodatek~\ref{app:locnica}.

\section{Razdelitev analize na več korakov}
\label{sec:koraki}

Celotno analizo bi lahko opravili z enim samim klicem modela. Model bi dobil prepis in eno dolgo navodilo, naj izlušči argumente, oceni njihovo
zgradbo, poišče odgovore in zmote ter napiše povzetek. Ta pristop smo zavrnili iz treh razlogov, odločitev pa je potrdil tudi preizkus.

Prvič, dolg prepis in več hkratnih nalog si delijo isto omejeno pozornost modela. Kakovost posameznih nalog zato pada, izhod pa rad preseže dovoljeno
dolžino in se odreže sredi stavka. Drugič, en klic pomeni en model, zato bi vse naloge opravil najdražji, čeprav bi za nekatere zadostoval cenejši. Tretjič, ena
napaka bi podrla celoto, pri razdeljeni analizi pa odpove le prizadeti korak, ostali se izvedejo do konca.

Pristop z enim klicem smo tudi preizkusili na posnetku, ki ni del zbirke za vrednotenje. Izidi so bili manj natančni kot pri razdelitvi na korake: model je vrnil manj virov, sestavljenih trditev ni razgradil in je vrnil več argumentov, kot jih je bilo v razpravi dejansko predstavljenih. 

Analizo smo zato razdelili na pet zaporednih korakov. Korak je stopnja analize, ki dobi določen vhod in vrne izhod, na katerem gradijo naslednje stopnje.
Štirje koraki so klici jezikovnega modela, in sicer izluščanje argumentov, zaznava zmot, preslikava odgovorov in sinteza. Vsak od njih opravi svojo nalogo
z enim samim pozivom. Tretji korak, preverjanje dejstev, je drugačen. Model v njem sicer sodeluje, vendar ne vrača presoje argumentacije, temveč izrečene podatke sooči z zunanjimi viri, in to z več klici na različne ponudnike, katerih izide nato združi sistem (razdelek~\ref{sec:fact-checking}). Vseh pet povzema
tabela~\ref{tbl:koraki}.
\begin{table}[H]
\centering
\small
\arrayrulecolor{crta}
\begin{tabular}{@{}p{0.28\textwidth}p{0.22\textwidth}p{0.42\textwidth}@{}}
\toprule
Korak & Glavni vhod & Izhod \\
\midrule
1 izluščanje argumentov & prepis & argumenti s premisami, stališča govorcev \\
\midrule
2 zaznava zmot & samo izhod 1 & zmote z imenom, kategorijo, delom argumenta, v katerem je napaka, in oznako argumenta \\
\midrule
3 preverjanje dejstev & premise iz izhoda 1 & razsodbe z navedenimi viri \\
\midrule
4 preslikava odgovorov & prepis, izhod 1 & zavrnitve, izmikanja \\
\midrule
5 sinteza & izhodi od 1 do 4 & povzetek, pri enem govorcu še seznam nepodprtih trditev \\
\bottomrule
\end{tabular}
\caption{Koraki analize z vhodi in izhodi.}
\label{tbl:koraki}
\end{table}

Razdelitev nam omogoča tudi, da za vsak korak izberemo model glede na zahtevnost naloge. Kateri model opravlja posamezni korak in zakaj, navaja dodatek~\ref{app:modeli}. Izluščanje je med koraki najpomembnejše, saj kot edino ustvari argumente, na katerih delajo vsi nadaljnji koraki. Odpoved izluščanja analizo ustavi, odpoved katerega koli drugega koraka pa izid le oslabi.

Iz tabele je razvidno, da vsak korak dobi le tisto, kar za svojo nalogo
potrebuje. Prepis prejmeta samo izluščanje argumentov in preslikava odgovorov. Ostali trije delajo z izhodi prejšnjih korakov. 
Zaznava zmot ne potrebuje celega prepisa, ker lahko samo iz izluščenih argumentov poišče zmote. Prejme argumente vseh govorcev, ne le enega.
Preverjanje dejstev dela na premisah in ne na prepisu, sicer bi preverjalo tudi podatke, ki v nobenem argumentu ne nastopajo, razsodb pa ne bi bilo mogoče pripeti k premisi.
Preslikava odgovorov je predvsem vezava že izluščenih argumentov na mesta v prepisu in sklepanja ne presoja, zato kot edini od štirih klicev analize teče na cenejšem modelu. Pri posnetku z enim govorcem odpade, ker ni sogovornika.

\chapter{Implementacija}
\label{ch:implementacija}

V tem poglavju so podrobno opisani posamezni koraki cevovoda. Zaledni del
sistema je napisan v jeziku Python, uporabniški vmesnik pa v ogrodju React. Potek od posnetka do poročila prikazujejo slike~\ref{fig:cevovod}--\ref{fig:cevovod3}.

\begin{figure}[H]
\centering
\includegraphics[width=0.82\textwidth]{loading.png}
\caption{Prikaz poteka analize. Zgoraj je opomba »Analysis in progress« in identifikator naloge. Sledi seznam korakov obdelave: zelene kljukice označujejo že dokončane korake, trenutni korak je označen z vijolično, preostali koraki so še v čakanju. Na dnu je napredovalni trak, ki prikazuje trenutno delo.}
\label{fig:loading}
\end{figure}

\begin{figure}[H]
\centering
\includegraphics[width=0.95\textwidth,keepaspectratio]{cevovod1.pdf}
\caption{Prvi del cevovoda, od vhoda do pripravljenega prepisa. Ob vsakem koraku je
navedeno, kaj v njem nastane, barva okvira pa pove, ali korak opravi zunanja
storitev, klic jezikovnega modela ali sam sistem.}
\label{fig:cevovod}
\end{figure}

\section{Pridobivanje zvoka in prepis}
\label{sec:prenos}

Sistem sprejme povezavo do posnetka na YouTubu ali datoteko, ki jo uporabnik
naloži sam. Pri povezavi pred prenosom poizve po metapodatkih in iz njih razbere naslov, ime kanala in trajanje. Uporabniku tako takoj pove, ali je posnetek dosegljiv, in mu ponudi izbiro odseka. Kadar trajanja ni mogoče pridobiti, mora uporabnik sam izbrati odsek, dolg največ 23 minut. Posnetek, daljši od 23 minut, in prenos v živo analiza zavrne. Meja je takšna, ker model v enem klicu sprejme največ 1400 sekund zvoka.

Prenos nato opravi orodje \texttt{yt-dlp}. To iz posnetka izlušči zvočno sled in jo shrani v formatu m4a. Slike ne prenese, ker je analiza govora ne potrebuje. Če je uporabnik izbral časovni odsek, orodje \texttt{ffmpeg} izreže samo izbrani del. Pred prepisom ffmpeg zvok po potrebi pretvori še v stisnjen zapis mp3, da ne preseže največje velikosti datoteke, ki jo sprejme model za prepis.

\subsection{Prepis z ločevanjem govorcev}
\label{ssec:diarizacija}

Prepis opravi model \texttt{gpt-4o-transcribe-diarize}, ki poleg besedila vrne tudi oznake govorcev.

Oznake, ki jih vrne prepis, so brezimenske, torej »Speaker~1«, »Speaker~2«. Kdo je kdo, pove uporabnik. Imena lahko vpiše pred analizo, vendar jih sistem vstavi šele po njej, po vrstnem redu govorcev: prvo ime dobi \sn{Speaker~1}, drugo \sn{Speaker~2}. Govorca lahko preimenuje tudi v vmesniku po analizi (razdelek~\ref{ssec:urejanje}).

Pri razpravi sistem pred izluščanjem preveri še, ali prepis vsebuje vsaj dva različna govorca. Če je ločevanje oba glasova zlilo v enega, analizo zavrne, saj razprave iz takega prepisa ni mogoče analizirati.

Prepis je shranjen kot datoteka .txt, v kateri je vsaka izjava v svoji vrstici v obliki \texttt{Ime: besedilo}. Najprej nastane v delovni mapi opravila \texttt{jobs/<oznaka opravila>/}, v kateri cevovod hrani preneseni zvok, vmesne datoteke in izide ter iz katere prepis bere tudi analiza. Nato ga strežnik prekopira v datoteko \texttt{transcripts/<oznaka opravila>.txt}, iz katere ga bere ponovna analiza (razdelek~\ref{ssec:rerun}). V podatkovni bazi prepis ni shranjen. Spodaj je odsek enega od shranjenih prepisov.

\begin{quote}
\footnotesize

\textbf{Dr.\ Raoul Kaninadyke:} Oh, disagree.

\textbf{Tristan Hughes:} Can I put this here? I don't know if this is going to be\ldots Can we do strongly disagree or strongly disagree?

\textbf{Dr.\ Raoul Kaninadyke:} You know, I'm only going to put disagree actually, because I think there's a way that you can\ldots

\textbf{Tristan Hughes:} There is, of course.
\end{quote}

\section{Večkoračna analiza argumentacije}
\label{sec:analiza}

\begin{figure}[H]
\centering
\includegraphics[width=0.95\textwidth,keepaspectratio]{cevovod2.pdf}
\caption{Drugi del cevovoda, večkoračna analiza argumentacije. Preverjanje dejstev je na sliki poenostavljeno na tri dele, podrobneje ga opisuje razdelek~\ref{sec:fact-checking}. }
\label{fig:cevovod2}
\end{figure}

Analiza je razdeljena na zaporedne korake, pri razpravi dveh govorcev na pet in pri nastopu enega govorca na štiri. Vsak dobi svoj določen vhod, opravi nalogo in vrne strukturiran izhod, ki ga uporabijo poznejši koraki. 

Prvi korak je izluščanje, ker vsi nadaljnji koraki delajo na argumentih, ki jih vrne. Zaznava zmot in preverjanje dejstev sta neodvisna drug od drugega, oba pa potrebujeta že izluščene argumente. Preslikava odgovorov potrebuje poleg argumentov še prepis, sinteza pa stoji zadnja, ker je edina, ki potrebuje izhod vseh prejšnjih. Kateri jezikovni model opravlja katero nalogo in zakaj je bil izbran, je navedeno v dodatku~\ref{app:modeli}. Pozivi, po katerih se ravnajo modeli, so v repozitoriju sistema v mapi \texttt{prompts}, vsak v svoji datoteki in razvrščeni po korakih (\url{https://github.com/dilan23Jill/Analiza-razprav}).

Pravila v pozivih smo določili sami na podlagi ogledov širokega nabora spletnih razprav, tako da so uporabna za razprave vseh vrst, na primer zgodovinske, filozofske ali politične. Osnutke pozivov smo nato z jezikovnim modelom preoblikovali v jasnejše besedilo, vsebine pravil pa pri tem nismo spreminjali.

\subsection{Sistemski pozivi}
\label{ssec:sistemski-poziv}

Vhod ni pri vseh korakih enak. Prepis posnetka dobita za vhod samo izluščanje argumentov in preslikava odgovorov. Pri teh je uporabniški poziv sestavljen iz prepisa, za katerim sledi navodilo. Prepis je v obeh klicih isti in na istem mestu, razlikuje ju le navodilo, ki mu sledi.

Zaznava zmot in sinteza prepisa ne prejmeta, ampak izluščene argumente in izhode prejšnjih korakov, ki so kot JSON vpisani v uporabniški poziv. 

Sistemski poziv razsodbe modelu pove le, da je preverjevalec dejstev in da mora odgovoriti z veljavnim JSON. Vse ostalo je v uporabniškem pozivu, skupaj s trditvijo in zbranim gradivom: merila, kdaj je trditev resnična, delno resnična, zavajajoča, neresnična ali nepreverljiva, prepoved, da bi model sam iskal po spletu ali se opiral na lastno znanje, ter navodilo, naj vsakemu viru posebej pripiše, kam kaže.

Klici, ki presojajo argumentacijo govorcev, dobijo v sistemskem pozivu isto skupno navodilo: poročajo nevtralno, ne razglasijo zmagovalca in ne ocenjujejo, čigav primer je močnejši. Ostali klici dobijo vsak svoje navodilo, zbiranje virov pa sistemskega poziva nima. Kaj vsebuje sistemski poziv posameznega klica, povzema tabela~\ref{tbl:hisna}.

\begin{table}[H]
\centering
\small
\arrayrulecolor{crta}
\begin{tabular}{@{}p{0.32\textwidth}p{0.13\textwidth}p{0.47\textwidth}@{}}
\toprule
Klic modela & Skupno navodilo & Kaj še vsebuje sistemski poziv \\
\midrule
izluščanje argumentov (1.~korak) & ne & naj argumente izlušči in ne presoja, argumenta ne sme izpustiti ali oslabiti, ter kdo v posnetku šteje za udeleženca \\
\midrule
zaznava zmot (2.~korak) & da & kdaj se zmota sme zapisati in kako se zapiše dvoumen primer \\
\midrule
izbira in razgradnja trditev (3.~korak) & ne & celotno navodilo za nalogo \\
\midrule
zbiranje virov (3.~korak) & ne & sistemskega poziva nima, navodilo je v uporabniškem pozivu \\
\midrule
razsodba (3.~korak) & ne & vlogo preverjevalca dejstev in zahtevo po veljavnem JSON \\
\midrule
preslikava odgovorov (4.~korak) & da & kaj šteje za izmikanje in kdo v posnetku šteje za udeleženca \\
\midrule
sinteza (5.~korak) & da & da sme poročati le o tem, kar so našli prejšnji koraki \\
\bottomrule
\end{tabular}
\caption{Vsebina sistemskega poziva posameznega klica modela.}
\label{tbl:hisna}
\end{table}

\subsection{Prvi korak: izluščanje argumentov}
\label{ssec:prehod1}

Prvi korak je edini, ki argumente ustvari. Vsi nadaljnji koraki delajo na njegovem izhodu, zato je ta korak najpomembnejši.

Da model pri istem prepisu ne vrača vsakič drugačnih argumentov, poziv presoje ne prepušča njemu, ampak določa pravilo: sklep se upošteva natanko takrat, ko ga govorec podpre z vsaj dvema premisama. Nič drugega ne odloča, ne kako pomemben se sklep zdi, ne koliko argumentov je model že vrnil.

Sklep z eno samo premiso ne postane samostojen argument. Tak sklep je skoraj vedno del širšega argumenta, zato ga model po navodilu tja pripne kot premiso. Koda pravilo po koraku še preveri in argument z manj kot dvema premisama odstrani, razen če bi govorec s tem ostal brez argumentov. Poziv tudi pove, da ciljnega števila argumentov ni. Če govorec ne sestavi nobenega celega argumenta, je pravilen izid prazen seznam.

Kadar več sklepov podpira isti širši sklep, poziv modelu naroča, naj jih ne vrne kot ločene argumente, ampak kot premise širšega sklepa. Tako nastane en argument. Poseben del poziva loči naštevanje argumentov od naštevanja premis: če govorec našteva korake ene same izpeljave, gre za en argument z več premisami in ne za več argumentov.

Primer iz nastopa o tehnologiji v šolah: govorka izpelje dva sklepa, vsakega iz lastnih premis: \sn{učenci si bolje zapomnijo besedilo, ki ga preberejo na papirju, kot tisto na zaslonu} in \sn{telefoni v razredu odvračajo pozornost od pouka}. Oba podpirata širši sklep \sn{tehnologija v razredu škoduje učenju}. Model iz tega vrne en argument s tem sklepom in dvema premisama, pri čemer vsaka premisa obdrži razlog, ki ga je zanjo navedla govorka.

Premisa pri tem ni gol podatek, ampak kratek podargument: trditev skupaj z razlogom, ki ga je zanjo navedel govorec. Več podatkov, ki podpirajo isti razlog, poziv združi v eno premiso, zato se število premis ne povečuje s številom navedenih statistik.

Sklep in premise niso dobesedni navedki iz prepisa. Poziv modelu naroča, naj jih povzame s svojimi besedami, pri tem ohrani vsebino, ki jo je izrekel govorec, izpusti pa ponavljanja, nedokončane povedi in primere, ki za sklepanje niso potrebni. Sklepa ne sme izpeljati sam in argumentu ne sme dodati razlogov, ki jih govorec ni navedel. Če govorec premise ne utemelji, se premisa zapiše kot gola trditev.

Prvi korak vrne tudi vlogo vsakega udeleženca: glavni govorec, debater ali moderator. Poziv modelu naroča, naj vlogo določi po tem, kako se udeleženec v pogovoru obnaša. Moderator skoraj izključno sprašuje, povzema in predaja besedo, lastnega stališča pa ne zagovarja. Poziv izrecno pove, da mnoge razprave moderatorja nimajo in da vloge ne sme vsiliti. Argumenti se izluščijo le za debaterje oziroma glavnega govorca. Moderator in drugi stranski glasovi, na primer občinstvo, se ne analizirajo. Vprašanja moderatorja model po navodilu uporabi kot kontekst: kratek odgovor debaterja, ki ima pomen le skupaj z vprašanjem, dopolni z vsebino vprašanja, da argument stoji sam zase. Za moderatorja korak zapiše le vprašanja, ki jih je zastavil, in katerega debaterja je bolj pritiskal, kar sinteza uporabi pri opisu poteka razprave.

Pri razpravi koda izid prvega koraka še preveri. Udeleženca, ki ga je model označil kot moderatorja, a ga kljub temu uvrstil med govorce, iz seznama odstrani. Če nato v seznamu ne ostaneta natanko dva govorca ali če model v metapodatkih označi več kot dva debaterja oziroma nobenega nasprotnika, sistem analizo ustavi, namesto da bi analiziral napačen par. Uporabnik dobi sporočilo z razlogom, pri enem samem govorcu pa napotek, naj analizo zažene v načinu za enega govorca.

Vhod koraka je prepis. Izhod ima dva dela. Metapodatki (\texttt{metadata}) vsebujejo temo, vlogo vsakega udeleženca, zapis o moderatorju ter zastavici, ali ima razprava več kot dva debaterja oziroma nobenega nasprotnika. Slovar govorcev (\texttt{speakers}) za vsakega govorca vsebuje stališče (\texttt{position}), torej v enem stavku povzeto, kar zagovarja, in seznam argumentov. Stališče argumentov ne omejuje: argument, ki ga govorec utemelji ob drugi temi, se izlušči kot samostojen argument. Skrajšan primer izhoda je spodaj. Oznaka \texttt{arg\_id} je v primeru že vpisana, čeprav jo koda doda šele po koraku.

\begin{lstlisting}[basicstyle=\ttfamily\scriptsize,breaklines=true]
{"metadata": {"topic": "Was the Trojan War historical?",
   "participants": {"Speaker A": "debater", "Speaker B": "debater",
                    "Speaker C": "moderator"},
   "moderator": {"present": true, "name": "Speaker C", "question_count": 7,
                 "questions": ["..."], "pressed_more": "Speaker A",
                 "notes": "..."},
   "too_many_debaters": false, "too_few_debaters": false},
 "speakers": {"Speaker A": {
   "position": "The Trojan War as Homer tells it did not happen.",
   "arguments": [{
      "arg_id": "Speaker A#0",
      "argument": "Homer's war is not a historical event: the site shows
                   destruction, not a ten-year siege.",
      "premises": [
        "Excavations at Hisarlik show several destruction layers - but none
         that fits a single long siege.",
        "The epics were written down centuries after the supposed events."]}]}}}
\end{lstlisting}

Poziv je v repozitoriju v mapi \texttt{prompts/1\_izluscanje}.

\subsection{Drugi korak: logične zmote}
\label{ssec:prehod2}

Drugi korak prejme izključno izhod prvega koraka in ne dobi prepisa. 
Zmote, kot so prehitra posplošitev, zamenjava sosledja z vzročnostjo in sklicevanje na avtoriteto, so razvidne iz izluščenega argumenta samega, zato prepis za to nalogo ni potreben.

Poziv modelu naroča, naj vrne samo zmote. Za argumente brez zmote ne vrne ničesar, saj sistem ne ocenjuje, kako dober je argument. Vsak zapis ima ime zmote, kategorijo, obrazložitev, oznako argumenta, na katerega se nanaša, in tisti del argumenta, v katerem je napaka: sklep ali premiso, prepisano iz izhoda prvega koraka. Navedek zato ni dobeseden iz prepisa, ampak iz izluščenega argumenta, saj korak prepisa ne prejme.

Ime pove, za katero zmoto gre, na primer prehitra posplošitev ali slamnati mož. Kategorija pove, v katero od treh skupin iz razdelka~\ref{sec:zmote} ime spada: formalne zmote, neformalne zmote ali šibko sklepanje. Prehitra posplošitev je tako ime, šibko sklepanje pa njena kategorija. V zapisu je ime v polju \texttt{type}, kategorija pa v polju \texttt{category}.

Vhod koraka je izhod prvega koraka v obliki JSON, torej argumenti vseh govorcev z oznakami. Izhod je seznam zmot:

\begin{lstlisting}[basicstyle=\ttfamily\scriptsize,breaklines=true]
{"fallacies": [{
   "arg_id": "Speaker B#2", "premise_index": 1, "speaker": "Speaker B",
   "type": "hasty_generalization",
   "category": "weak_reasoning",
   "evidence": "Achilles was a terrible person - ...",
   "explanation": "One character is taken as proof of how ..."}]}
\end{lstlisting}

Ime zmote mora model po navodilu izbrati iz zaprtega slovarja 32 poimenovanih zmot, ki ga dopolnjuje še možnost \texttt{other} za primer, ko nobeno ime ne ustreza. Brez te omejitve isti pojav dobi vsakič drugo ime in ga med zagoni ni mogoče prešteti. Če model vrne ime zunaj te množice, ga shema poskusi preslikati na sopomenko, sicer pa ga ohrani tako, kot ga je zapisal model.

Kategorijo model sicer navede, vendar pri vseh 32 imenih iz slovarja obvelja tista, ki jo sistem izpelje iz imena. Potrjevanje posledice je vedno formalna zmota, slamnati mož vedno neformalna, prehitra posplošitev vedno šibko sklepanje. Uvrstitev v skupino je z imenom že določena, zato jo namesto modela opravi sistem. Modelova kategorija obvelja le pri imenu zunaj slovarja in pri \texttt{other}, kategorijo, ki je ne prepozna, pa shema nadomesti z neformalno.

Spodaj je kratek izsek iz poziva:

\begin{lstlisting}[basicstyle=\ttfamily\scriptsize,breaklines=true]
PART B - RHETORIC != FALLACY - CLASSIFY CORRECTLY:
These are rhetorical DEVICES. Do NOT report them as fallacies when they accompany
real argumentation:
  * Hyperbole and dramatic emphasis        * Rhetorical questions
  * Analogy, metaphor, vivid imagery       * Personal anecdote used as illustration
  * Humor, irony, sarcasm                  * Anaphora / repetition for emphasis
  * Framing and loaded word choice         * Appeals to shared values
The SAME move becomes a fallacy ONLY when it REPLACES the argument (e.g. emotional
appeal with no reasoning = appeal to emotion; anecdote presented as proof of a
general rule = hasty generalization). Ask: "If I strip this device away, is there
still an argument left?" Yes -> rhetoric, do not report it. No -> consider a fallacy.
[...]
AMBIGUOUS CASES - IMPORTANT:
Sometimes a statement COULD be a fallacy OR a legitimate rhetorical device - it
depends on interpretation. [...] In these cases:
- Still include it, but say so in the explanation
- Present BOTH interpretations: "This could be seen as [fallacy] because [...],
  but it could also be interpreted as [legitimate use] because [...]."
- Let the reader decide - your job is to flag it and explain both sides.
\end{lstlisting}

Dvoumnega primera torej korak ne izpusti, ampak ga zapiše in v obrazložitvi vrne dve branji. Odločitev je prepuščena bralcu, ki zaznavo v aplikaciji lahko potrdi, zavrne ali popravi, lahko pa doda tudi zmoto, ki je model ni vrnil (razdelek~\ref{ssec:urejanje}).

\subsection{Tretji korak: preverjanje dejstev}
\label{sec:fact-checking}

Ta korak iz premis izbere podatke, ki jih je mogoče preveriti z zunanjimi viri, zanje zbere gradivo in o vsakem vrne razsodbo. Za razliko od drugih korakov ni en klic modela, ampak zaporedje stopenj, opisanih v nadaljevanju: izbira trditev, razgradnja sestavljenih trditev, usmerjanje k zbiralcem in zbiranje gradiva, sestava seznama virov in razsojanje.

V tem koraku se uporabljajo štirje izrazi:
\begin{itemize}
\item \textbf{zbirka} je zunanja podatkovna zbirka ali storitev, v kateri sistem išče, na primer PubMed ali Wikidata,
\item \textbf{zbiralec} je del sistema, ki za trditev poizve v eni ali več zbirkah in vrne, kar najde,
\item \textbf{gradivo} je vse, kar zbiralci vrnejo za eno trditev,
\item \textbf{vir} je posamezen članek, stran ali zapis iz gradiva, ki ga sistem uvrsti na seznam virov ob razsodbi.
\end{itemize}

Vhod koraka je izhod prvega koraka. Koda ga modelu ne pošlje kot JSON, ampak ga pretvori v besedilo, po en blok za vsak argument:

\begin{lstlisting}[basicstyle=\ttfamily\scriptsize,breaklines=true]
[Speaker B#2] speaker: Speaker B
  position: Big cities carry most of the country's weight: ...
  premise 0: France has about 68 million inhabitants, and more than two
             million of them live in Paris.
  premise 1: Cities are more pleasant to live in than the countryside.
\end{lstlisting}

Vrstica \texttt{position} vsebuje besedilo argumenta, torej sklep z jedrom utemeljitve, in ne stališča govorca. Izhod koraka je za vsako trditev en zapis z razsodbo in viri, ki je prikazan na koncu razdelka~\ref{ssec:razsojanje}.

\subsubsection{Izbira preverljivih trditev}
\label{ssec:claim-extraction-fc}

Ta stopnja dobi izluščene argumente, tako kot drugi korak, in iz njihovih premis izbere tiste, ki so preverljive. Nalogo opravi model \texttt{gpt-5.6-luna} po navodilu, ki je v dodatku~\ref{app:izbira-trditev}.

Navodilo modelu naroča, naj loči dejstveno trditev od zagovarjanega mnenja. Premise, ki izražajo vrednote, se ne preverjajo. Navodilo prepoveduje tudi posploševanje: model mora vsako število, datum in besedo, ki določa obseg, prenesti natanko tako, kot je v premisi. Iz premise \sn{brezposelnost v Sloveniji se je leta 2023 znižala za dve odstotni točki} mora nastati trditev z vsemi tremi podatki, ne pa na primer \sn{brezposelnost se znižuje}.

V istem odgovoru model po navodilu vsaki trditvi pripiše tip, kar je vrstica \texttt{claim\_type} v navodilu (dodatek~\ref{app:izbira-trditev}). Tipe smo izbrali sami na podlagi poznavanja različnih vrst razprav. Trditev o številu obolelih je \texttt{health}, trditev o letnici bitke \texttt{historical}, dobesedno povzemanje tujih besed \texttt{quote}. Glede na dodeljeni tip sistem izbere zbiralce (tabela~\ref{tbl:usmerjanje}). Tip, ki ni del množice, sistem preslika na \texttt{statistic}, pri katerem stečejo vsi zbiralci. Zgornje meje števila preverjenih trditev sistem nima.

Izhod je seznam \texttt{claims}, v katerem ima vsaka trditev besedilo, oznako argumenta, številko premise, govorca, tip in en stavek o tem, čemu jo argument rabi. Iz premise \sn{Francija ima okoli 68 milijonov prebivalcev, v Parizu pa jih živi več kot dva milijona} zgornjega argumenta tako nastaneta dve trditvi, premisa \sn{mesta so za življenje prijetnejša od podeželja} pa ne da nobene, ker izraža mnenje:

\begin{lstlisting}[basicstyle=\ttfamily\scriptsize,breaklines=true]
{"claims": [
  {"exact_claim": "France has about 68 million inhabitants",
   "arg_id": "Speaker B#2", "premise_index": 0, "speaker": "Speaker B",
   "claim_type": "statistic",
   "context": "Used to show how large the country's population is."},
  {"exact_claim": "More than two million people live in Paris",
   "arg_id": "Speaker B#2", "premise_index": 0, "speaker": "Speaker B",
   "claim_type": "statistic",
   "context": "Used to show how much of the population lives in the capital."}]}
\end{lstlisting}

\subsubsection{Razgradnja sestavljenih trditev}
\label{ssec:clustering}

Govorec lahko v enem stavku pove dve različni dejstvi. Če je eno točno in drugo ne, razsodba za tak stavek ni mogoča, saj bi morala biti hkrati resnična in neresnična. Razgradnja zato tako trditev razstavi na več podtrditev, ki se preverijo ločeno. Vse trditve model prejme v enem klicu kot oštevilčen seznam in po navodilu razstavi le tiste, ki vsebujejo več neodvisnih dejstev, ostale pa vrne nespremenjene. Vsaka podtrditev obdrži govorca, oznako argumenta in številko premise, iz katere je nastala. Pri preizkusih je za to nalogo zadoščal najcenejši model (\texttt{gpt-4.1-nano}). Navodilo je v repozitoriju v datoteki \texttt{prompts/3\_preverjanje\_dejstev/razgradnja.txt}.

Vhod je oštevilčen seznam trditev, vsaka s svojim tipom:

\begin{lstlisting}[basicstyle=\ttfamily\scriptsize,breaklines=true]
1. [economic] "Unemployment fell by two percentage points in 2023 and
   inflation fell below three percent"
2. [historical] "The epics were written down centuries after the events"
\end{lstlisting}

Prva trditev vsebuje dve neodvisni dejstvi, zato iz nje nastaneta dve podtrditvi, druga pa ostane nespremenjena. Izhod je seznam \texttt{decomposed}, v katerem je vsaka podtrditev svoj zapis. Podtrditve iste izvorne trditve imajo v polju \texttt{parent\_claim} njeno besedilo, trditve z enim samim dejstvom pa model vrne nespremenjene in brez tega polja. Govorca, oznako argumenta in številko premise podtrditvam pripiše koda iz izvorne trditve, zato jih model ne vrača. Izhod za zgornji vhod:

\begin{lstlisting}[basicstyle=\ttfamily\scriptsize,breaklines=true]
{"decomposed": [
  {"exact_claim": "Unemployment fell by two percentage points in 2023",
   "claim_type": "economic",
   "parent_claim": "Unemployment fell by two percentage points in 2023 and ..."},
  {"exact_claim": "Inflation fell below three percent in 2023",
   "claim_type": "economic",
   "parent_claim": "Unemployment fell by two percentage points in 2023 and ..."},
  {"exact_claim": "The epics were written down centuries after the events",
   "claim_type": "historical"}]}
\end{lstlisting}

\subsubsection{Usmerjanje poizvedb k zbiralcem}
\label{ssec:engine-routing}

Sistem ima šest zbiralcev: akademski zbiralec, ki hkrati poizveduje v zbirkah znanstvenih člankov PubMed, Semantic Scholar, OpenAlex in Crossref, zbiralec Wikidata, ki poizveduje v zbirki strukturiranih podatkov Wikidata, zbiralec Google Fact Check, ki poizveduje v zbirki objavljenih preverjanj, ter tri zbiralce, ki uporabljajo jezikovne modele z dostopom do spleta: spletno iskanje, Perplexity in Grok. Nobeden ne vrne razsodbe, ampak le gradivo, razsodbo pa nato vrne en sam klic modela (razdelek~\ref{ssec:razsojanje}).

Poizvedbe se ne pošljejo vsem zbiralcem. Katere zbiralce trditev sproži, določa njen tip, saj ta pove, kje jo je smiselno iskati. Razporeditev, ki smo jo določili sami, navaja tabela~\ref{tbl:usmerjanje}, razloge zanjo pa odstavek pod tabelo.

\begin{table}[H]
\centering
\small
\arrayrulecolor{crta}
\begin{tabular}{@{}llcccccc@{}}
\toprule
Tip trditve & oznaka & \rot{akademske zbirke} & \rot{Wikidata} &
\rot{Google Fact Check} & \rot{Perplexity} & \rot{Grok} & \rot{spletno iskanje} \\
\midrule
statistična & \texttt{statistic} & $\bullet$ & $\bullet$ & $\bullet$ & $\bullet$ & $\bullet$ & $\bullet$ \\
politična & \texttt{policy} &  & $\bullet$ & $\bullet$ & $\bullet$ & $\bullet$ & $\bullet$ \\
ekonomska & \texttt{economic} &  & $\bullet$ & $\bullet$ & $\bullet$ & $\bullet$ & $\bullet$ \\
navedek & \texttt{quote} &  & $\circ$ & $\bullet$ & $\bullet$ & $\bullet$ & $\bullet$ \\
znanstvena & \texttt{scientific} & $\bullet$ & $\circ$ & $\bullet$ & $\bullet$ &  & $\bullet$ \\
zdravstvena & \texttt{health} & $\bullet$ & $\circ$ & $\bullet$ & $\bullet$ &  & $\bullet$ \\
zgodovinska & \texttt{historical} &  & $\bullet$ & $\bullet$ & $\bullet$ &  & $\bullet$ \\
zemljepisna & \texttt{geographic} &  & $\bullet$ & $\bullet$ & $\bullet$ &  & $\bullet$ \\
\bottomrule
\end{tabular}
\caption{Zbiralci, ki stečejo pri posameznem tipu trditve. Polna pika
$\bullet$ pomeni, da zbiralec steče vedno, prazni znak $\circ$, da steče le, kadar trditev vsebuje letnico ali katero koli število, prazno polje pa, da pri tem tipu ne steče. Stolpec akademske zbirke je akademski zbiralec, ki poizveduje v zbirkah PubMed, Semantic Scholar, OpenAlex in Crossref. Google Fact Check, Perplexity in spletno iskanje stečejo pri vseh tipih.}
\label{tbl:usmerjanje}
\end{table}

Akademske zbirke se uporabijo pri znanstvenih, zdravstvenih in statističnih trditvah, ker spletno iskanje recenziranih člankov ne najde zanesljivo. Grok se pri znanstvenih, zdravstvenih, zgodovinskih in zemljepisnih trditvah izpusti, ker išče tudi po omrežju X, kjer je o dolgo uveljavljenih dejstvih več razprave kot gradiva. Wikidata je zbirka strukturiranih podatkov, zato se uporabi pri tipih, ki običajno vsebujejo merljive vrednosti, pri ostalih pa le, kadar je število v trditvi sami.

Zbiralci delujejo na dva načina. Trije zbiralci, akademski, Wikidata in Google Fact Check, kličejo spletne vmesnike (API) šestih zbirk neposredno s standardno knjižnico Python \texttt{urllib}. Ti vrnejo strukturirane zapise v obliki JSON ali XML, iz katerih sistem prebere posamezna polja. Trije uporabljajo jezikovne modele z dostopom do spleta, ki jih koda kliče s knjižnico \texttt{openai}. Ti vrnejo besedilo, ki ga zapiše model, in seznam povezav. Zbiralce prve vrste povzema tabela~\ref{tbl:zadetki}, druge vrste pa tabela~\ref{tbl:zbiralci-modeli}. Zadetki štirih akademskih zbirk se najprej združijo v en seznam, iz njega odpadejo enaki naslovi, preostali pa se uredijo po številu citatov in letnici.

Načina se razlikujeta tudi po tem, kaj zbiralec pošlje zbirki. Zbiralci prve vrste ne pošljejo trditve, ampak poizvedbo iz ključnih besed. Če trditev ni v angleščini, jo najcenejši model prevede v tri do šest angleških ključnih besed, sicer koda iz trditve odstrani pogoste besede in obdrži prvih šest. Wikidata prejme le lastna imena iz trditve, največ tri, in jih poišče kot entitete. Zbiralci druge vrste pa modelu pošljejo celotno besedilo trditve skupaj s kontekstom, poizvedbe po spletu pa sestavi model sam. Pri prvih je izid odvisen od tega, katere besede so v poizvedbi, pri drugih pa od tega, kako model zastavi iskanje.

\begin{table}[H]
\centering
\small
\arrayrulecolor{crta}
\begin{tabular}{@{}p{0.18\textwidth}p{0.30\textwidth}p{0.13\textwidth}p{0.31\textwidth}@{}}
\toprule
Zbirka & Klic & Največ zadetkov & Kaj vrne \\
\midrule
PubMed & NCBI E-utilities (\texttt{esearch}, \texttt{efetch}), XML & pet & naslov, revija, letnica, izvleček do 600 znakov \\
\midrule
Semantic Scholar & Graph API, JSON & pet & naslov, revija, letnica, število citatov, izvleček \\
\midrule
OpenAlex & REST API, JSON & pet & naslov, revija, letnica, število citatov, izvleček \\
\midrule
Crossref & REST API, JSON & pet & naslov, revija, letnica, število citatov, avtorji, izvleček \\
\midrule
Wikidata & iskanje entitet (\texttt{wbsearchentities}) in poizvedba SPARQL, JSON & tri & ime entitete, opis in nekaj lastnosti \\
\midrule
Google Fact Check & Fact Check Tools API (\texttt{claims:search}), JSON & tri & kdo je preverjal, njegova ocena, naslov in datum \\
\bottomrule
\end{tabular}
\caption{Zbirke, ki jih zbiralci kličejo prek spletnih vmesnikov. Iz štirih akademskih zbirk skupaj ostane največ šest člankov.}
\label{tbl:zadetki}
\end{table}

\begin{table}[H]
\centering
\small
\arrayrulecolor{crta}
\begin{tabular}{@{}p{0.18\textwidth}p{0.36\textwidth}p{0.38\textwidth}@{}}
\toprule
Zbiralec & Klic & Kaj vrne \\
\midrule
spletno iskanje & OpenAI Responses API, model \texttt{gpt-4o} z orodjem \texttt{web\_search\_preview} & JSON s povzetkom najdenega in seznamom strani z naslovom, povezavo, datumom in navedkom \\
\midrule
Perplexity & API Perplexity, model \texttt{sonar-pro}, po pet trditev v enem klicu & besedilo z oštevilčenim razdelkom za vsako trditev in seznam povezav, na katere se besedilo sklicuje \\
\midrule
Grok & API xAI, model \texttt{grok-4.3} & JSON s povzetkom in seznamom strani z naslovom, povezavo in navedkom \\
\bottomrule
\end{tabular}
\caption{Zbiralci, ki uporabljajo jezikovne modele z dostopom do spleta. Števila zadetkov ne omejujejo, zato je seznam virov omejen posebej.}
\label{tbl:zbiralci-modeli}
\end{table}

Zbiralec Perplexity trditve pošlje v paketih po pet in dobi oštevilčen seznam odgovorov, ki ga koda razreže po trditvah. Povezave Perplexity vrne za ves paket skupaj, zato trditvi pripadajo le tiste, na katere se njen razdelek sklicuje z oznako $[n]$. Brez tega bi seznam virov ene trditve vseboval strani o štirih drugih. Kadar odgovora ni mogoče razrezati, ga koda zavrže, namesto da bi gradivo pripisala napačni trditvi.

Vhod te stopnje je trditev iz prejšnjih stopenj: besedilo, tip in stavek o tem, čemu jo argument rabi. Izhod je gradivo, po en sklop za vsak zbiralec, in seznam zbiralcev, ki niso stekli ali niso vrnili ničesar, z razlogom (\texttt{skipped}). Za trditev \sn{France has about 68 million inhabitants}, ki je tipa \texttt{statistic}, stečejo vsi zbiralci. Skrajšan primer gradiva:

\begin{lstlisting}[basicstyle=\ttfamily\scriptsize,breaklines=true]
{"papers": [{"title": "...", "journal": "...", "year": 2022,
             "citations": 41, "abstract": "..."}],
 "wikidata": [{"entity": "France", "entity_id": "Q142",
               "fact_summary": "population: ..."}],
 "factchecks": [],
 "web": {"findings": "INSEE reports about 68.4 million inhabitants
                      on 1 January 2024 ...",
         "sources": [{"title": "...", "url": "https://www.insee.fr/...",
                      "date": "2024", "relevant_quote": "...",
                      "source_type": "official_stat"}]},
 "grok": {"findings": "...", "sources": [...]},
 "perplexity": null,
 "skipped": {"perplexity": "no answer, or the batch answer could not
                            be split per claim"}}
\end{lstlisting}

\subsubsection{Sestava seznama virov}
\label{ssec:viri}

Iz vsega gradiva koda sestavi en seznam z največ desetimi viri. Nobeden od zbiralcev ne pove, ali trditev drži.

Mesta se polnijo izmenično, po en vir od vsakega zbiralca, ki je kaj našel, nato drugi krog, dokler mest ni deset ali dokler zbiralcem ne zmanjka gradiva. Vrstni red kroga je stalen:

\begin{enumerate}\itemsep2pt
\item stran, ki jo je našlo spletno iskanje,
\item članek iz akademskih zbirk,
\item strukturirano dejstvo iz Wikidate,
\item objavljeno preverjanje iz Google Fact Check,
\item stran iz Grokovega iskanja po spletu in omrežju X,
\item povezava, na katero se sklicuje odgovor Perplexityja.
\end{enumerate}

Vsak zbiralec ponudi svoje zadetke v vrstnem redu, v katerem jih je uredil. Akademski jih uredi po številu citatov, zato v prvi krog vstopi najbolje citiran članek.

Izmenično polnjenje ohrani raznolikost virov. Deset člankov iste založbe pokrije eno ali dve domeni, dva članka, dve novici, podatek statističnega urada in objavljeno preverjanje pa več. Koda o virih prešteje prav dvoje, koliko jih je in s koliko različnih domen prihajajo, obe števili pa se izpišeta ob razsodbi, zato mora seznam to raznolikost ohraniti.

Vsak vir ima naslov in povezavo, kadar ju zbiralec vrne, pa še datum in kratek navedek iz besedila. Vir iz Perplexityja je le povezava, zato je njegov naslov kar ime domene. Vsaka povezava se uvrsti le enkrat. Kadar isto stran vrneta dva zbiralca, obvelja zapis tistega, ki je bil prej na vrsti, drugi pa svojega mesta ne izgubi, saj v istem krogu ponudi naslednjo stran. Ker seznam nastane izključno iz gradiva zbiralcev, razsojanje ne more dodati vira, ki ga sistem nima. Sistem virov tudi ne razvršča po kakovosti, saj to, kaj posamezen vir pove o trditvi, presodi razsojanje.

\label{ssec:neodvisnost}

Sistem neodvisnosti virov ne ugotavlja. Prešteje le, s koliko različnih domen prihajajo viri, in to število izpiše ob trditvi, saj pet zadetkov z istega spletišča ni pet potrditev. Povezavi \texttt{who.int/news/a} in \texttt{who.int/fact-sheets/b} sta dva vira in ena domena. Domena je ime gostitelja brez predpone \texttt{www.}, zato \texttt{data.who.int} in \texttt{who.int} štejeta za dve domeni. Viri z različnih domen so lahko kljub temu odvisni, na primer kadar več novic povzema isto poročilo, zato število domen ni mera neodvisnosti.

Števili virov in domen razsodbe ne spreminjata in v noben seštevek ne štejeta. Bralcu omogočata, da razsodbo, ki stoji na dveh virih ene domene, prebere drugače kot tisto, ki stoji na osmih virih iz petih domen.

Vhod te stopnje je gradivo iz prejšnje stopnje. Izhod je seznam virov, v katerem ima vsak vir enaka polja ne glede na to, kateri zbiralec ga je našel, ter število virov in število različnih domen (\texttt{evidence\_metrics}). Oba sta vidna v končnem zapisu na koncu razdelka~\ref{ssec:razsojanje}.

\subsubsection{Ena razsodba nad vsem zbranim gradivom}
\label{ssec:razsojanje}

Zbiranje in razsojanje sta ločeni stopnji. Ko so vsi izbrani zbiralci končali, gre njihovo gradivo v en sam klic modela, ki edini vrne razsodbo. Poziv modelu prepoveduje iskanje po spletu in rabo lastnega znanja, zato razsodbo izpelje le iz prejetega gradiva.

Poziv je v repozitoriju v datoteki \texttt{prompts/3\_preverjanje\_dejstev/razsodba.txt}, njegovo zgradbo pa prikazuje primer na koncu razdelka.

Model vrne razsodbo, obrazložitev in, kadar trditev ne drži ali je zavajajoča, v enem stavku še pravilen podatek.

V istem odgovoru mora model po navodilu vsakemu viru z oštevilčenega seznama pripisati eno od istih petih razsodb, tokrat za ožje vprašanje, in sicer kam kaže ta vir sam zase, sodeč po naslovu, datumu in navedku, ki ga je prejel. Vir, ki govori o čem drugem ali pove premalo, dobi \sn{nepreverljivo}, in poziv izrecno zahteva prav to vrednost povsod, kjer vir ne podpira nobene druge. Kako je to videti v vmesniku, prikazuje slika~\ref{fig:viri-okno}.

\begin{figure}[htbp]
\centering
\includegraphics[width=0.95\textwidth,keepaspectratio]{viri_okno.png}
\caption{Okno ene preverjene trditve. Zgoraj sta razsodba in obrazložitev, pod
njima seštevek oznak virov, spodaj pa oštevilčen seznam desetih virov, kjer ima
vsak svojo oznako. Seštevek oznak je le izpis. Razsodba iz njega ni izračunana, zato je trditev lahko \sn{resnično}, čeprav štirje viri o njej ne povedo ničesar.}
\label{fig:viri-okno}
\end{figure}

Seštevek oznak ne določa razsodbe, zato se razsodba od večine oznak sme razlikovati. Ena uradna statistika lahko odtehta tri strani, ki druga drugo prepisujejo, in poziv v tem primeru zahteva, da model razlog zapiše v obrazložitev. Kadar si gradivo nasprotuje, se to prav tako zapiše v obrazložitev, namesto da bi obveljal eden od virov.

Koda oznake preveri, preden jih sprejme. Številka, ki kaže mimo seznama, se zavrže, vrednost, ki ni ena od petih, postane \sn{nepreverljivo}, vir, ki ga model ni omenil, pa ostane brez oznake in se v seštevek ne šteje. Model tako ne more označiti vira, ki ga sistem nima, in tudi ne več virov, kot jih je dobil.

Vhod te stopnje je poziv, ki ga koda sestavi iz trditve, gradiva in seznama virov. Skrajšan primer za trditev o Franciji:

\begin{lstlisting}[basicstyle=\ttfamily\scriptsize,breaklines=true]
CLAIM: "France has about 68 million inhabitants"
TYPE: statistic
CLAIM CONTEXT: Stated while arguing: Used to show how large the
country's population is.

=== PEER-REVIEWED LITERATURE ===
1. ... (2022) - ... [41 citations]
=== STRUCTURED KNOWLEDGE (Wikidata) ===
- France (Q142): population: ...
=== WEB SEARCH FINDINGS ===
INSEE reports about 68.4 million inhabitants on 1 January 2024 ...
=== NUMBERED SOURCE LIST ===
[1] ... (official_stat, 2024) - https://www.insee.fr/...
[2] ... (peer_reviewed, 2022) - https://doi.org/...
[3] Wikidata: France (Q142) (structured_knowledge, Unknown) - ...
=== COLLECTORS THAT DID NOT RUN ===
- perplexity: no answer, or the batch answer could not be split per claim
\end{lstlisting}

\noindent Model vrne razsodbo, obrazložitev, oznako vsakega vira po številki in popravek:

\begin{lstlisting}[basicstyle=\ttfamily\scriptsize,breaklines=true]
{"verdict": "TRUE",
 "explanation": "INSEE gives about 68.4 million inhabitants ...",
 "sources": [{"n": 1, "verdict": "TRUE"}, {"n": 2, "verdict": "UNVERIFIABLE"},
             {"n": 3, "verdict": "TRUE"}, ...],
 "correction": ""}
\end{lstlisting}

Koda iz tega sestavi končni zapis trditve, ki je izhod celotnega tretjega koraka. Oznako vsakega vira vpiše k viru (\texttt{source\_verdict}), oznake sešteje (\texttt{source\_verdicts}) in doda število virov in število različnih domen. Vrednosti TRUE, PARTIALLY\_TRUE, MISLEADING, FALSE in UNVERIFIABLE ustrezajo razsodbam \sn{resnično}, \sn{delno resnično}, \sn{zavajajoče}, \sn{neresnično} in \sn{nepreverljivo}.

\begin{lstlisting}[basicstyle=\ttfamily\scriptsize,breaklines=true]
{"exact_claim": "France has about 68 million inhabitants",
 "arg_id": "Speaker B#2", "premise_index": 0, "speaker": "Speaker B",
 "claim_type": "statistic", "verdict": "TRUE",
 "explanation": "...", "correction": "",
 "sources": [{"title": "...", "url": "https://www.insee.fr/...",
              "date": "2024", "relevant_quote": "...",
              "source_type": "official_stat", "source_verdict": "TRUE"}, ...],
 "source_verdicts": {"TRUE": 4, "PARTIALLY_TRUE": 1, "MISLEADING": 0,
                     "FALSE": 0, "UNVERIFIABLE": 5},
 "evidence_metrics": {"source_count": 10, "independent_domain_count": 8}}
\end{lstlisting}

\subsection{Četrti korak: preslikava odgovorov in izmikanj}
\label{ssec:prehod4}

Ta korak poveže odgovore z argumenti, na katere se nanašajo, in zabeleži, na katera vprašanja govorci niso odgovorili. Za vsak odgovor navede, kdo je komu
odgovoril, na kateri argument, jedro odgovora in odziv izvirnega govorca v enem stavku, ter vrsto odgovora. Vrste so štiri, in sicer neposredno nasprotovanje,
rušenje premise, druga razlaga in izpodbijanje sklepanja.

Za vsako izmikanje zapiše, kdo se je izmaknil, na katero vprašanje, na kakšen način in kolikokrat je bilo vprašanje postavljeno. Načini so preusmeritev,
menjava teme, odgovor brez vsebine, delni odgovor in preglasitev. Obe množici sta zaprti. Vrednost zunaj zaprte množice shema preslika na sopomenko, sicer pa na privzeto vrednost
(razdelek~\ref{ssec:validacija}).

Vhod koraka sta prepis kot besedilo in izhod prvega koraka kot JSON z oznakami argumentov, izhod pa sta seznama zavrnitev in izmikanj:

\begin{lstlisting}[basicstyle=\ttfamily\scriptsize,breaklines=true]
{"rebuttals": [{"by": "Speaker B", "to": "Speaker A",
   "target_arg_id": "Speaker A#0",
   "rebuttal_type": "alternative_explanation",
   "rebuttal_content": "The layers fit a long conflict remembered as one war.",
   "response": "Speaker A concedes a conflict but not Homer's version."}],
 "evasions": [{"evading_speaker": "Speaker A",
   "question_asked": "Which layer would you accept as Troy?",
   "evasion_type": "partial_answer", "times_asked": 2,
   "explanation": "..."}]}
\end{lstlisting}

Kdo je iz izmenjave izšel bolje, korak ne vrne, ker mu poziv to prepoveduje. Zapiše, kaj je bilo rečeno,
presojo pa prepusti bralcu. Tudi izmikanje ni stopnjevano: sistem zapiše, da
odgovora ni bilo, na kakšen način in kolikokrat, ne pa, kako hudo je to. Prav to zahteva tudi poziv:

\begin{lstlisting}[basicstyle=\ttfamily\scriptsize,breaklines=true]
Record what was said, not who you think won. Whether the argument survived the
exchange is left to the reader.
\end{lstlisting}

Ta korak je poleg prvega edini, ki prejme prepis, ker potrebuje cel kontekst za iskanje odgovorov ali izmikanj. V navodilu ima pravila o moderatorju in o glasovih iz občinstva: vprašanje moderatorja ni zavrnitev, izmikanje pa šteje le med govorcema. Izvede se le v načinu razprave, pri posnetkih z enim govorcem pa se izpusti. Navodilo je v repozitoriju v mapi \texttt{prompts/4\_zavrnitve}.

\subsection{Peti korak: sinteza}
\label{ssec:prehod5}

Ta korak prejme izhode prvega, drugega in četrtega koraka kot JSON ter razsodbe tretjega kot besedilni seznam, prepisa pa ne. Vsak vhod je omejen na dolžino, ki jo določa koda; kadar ga je
treba skrajšati, se odstranjujejo zadnji elementi najdaljšega seznama, tako da
krajšanje ne odreže v celoti enega od govorcev.

Korak vrne povzetek razprave v štirih do šestih stavkih, ki opiše glavna
razhajanja, kaj je kdo zagovarjal in na čem je to utemeljil, vzorce izmikanja in
vprašanja, ki so ostala odprta. Za vsakega debaterja vrne še stavek o tem, kako je argumentiral (\texttt{rhetorical\_style}), in stavek, ki povzame razsodbe njegovih trditev (\texttt{factual\_accuracy}). V polje \texttt{moderator\_influence} zapiše, kako so vprašanja moderatorja oblikovala izmenjavo.

\begin{lstlisting}[basicstyle=\ttfamily\scriptsize,breaklines=true]
{"comparative_evaluation": {
   "moderator_influence": {"present": false},
   "per_speaker": {"Speaker A": {
      "rhetorical_style": "Builds each point from an archaeological detail.",
      "factual_accuracy": "Most checked claims were rated TRUE."}}},
 "summary": "The main clash was over ..."}
\end{lstlisting}

Poziv modelu prepoveduje, da bi vrnil razsodbo o tem, čigav nastop je bil boljši, ali naštel močne in šibke točke. Utemeljitev te odločitve je v
razdelku~\ref{sec:kategorije-vs-stevilke}, poziv pa jo zapiše kot pravilo:

\begin{lstlisting}[basicstyle=\ttfamily\scriptsize,breaklines=true]
NO VERDICT: Do NOT declare a winner and do NOT rank the debaters. Never state
that one side "won", "dominated", "prevailed" or "made the stronger case", and
never award a category to either debater. Describe what each side argued and
how they argued it, and leave the conclusion to the reader. Do NOT rate the
quality of anyone's case. This applies to EVERY field below, including
`summary`.
\end{lstlisting}

Ker korak prejme izhode vseh prejšnjih, je tudi edino mesto, kjer se
ugotovitve o dejstvih in ugotovitve o sklepanju srečajo v istem besedilu. Pri
posnetku z enim govorcem dobi svoj poziv brez zavrnitev in brez primerjave med
govorcema, namesto tega našteje trditve, ki jih je govorec izrekel brez vsakega
razloga. Ker sinteza prepisa ne prejme, prvi korak pa gole trditve izpusti, lahko ta
seznam nastane le iz premis, ki jih govorec ni utemeljil. Trditev, ki ni postala
del nobenega argumenta, vanj ne pride. Obe navodili sta v repozitoriju v mapi \texttt{prompts/5\_sinteza}.

\section{Obdelava v kodi}
\label{sec:determinizem}

\begin{figure}[H]
\centering
\includegraphics[width=0.95\textwidth,keepaspectratio]{cevovod3.pdf}
\caption{Tretji del cevovoda, od obdelave izhodov do prikaza. Črtkana pot iz shrambe vodi na prvi korak analize (slika~\ref{fig:cevovod2}).}
\label{fig:cevovod3}
\end{figure}

Model se lahko zmoti v obliki odgovora ali v oznakah, ki jih vrne. Zato koda
vsak odgovor preveri in ga po potrebi popravi, preden gre naprej. 

\label{ssec:validacija}

Vsak odgovor koraka analize gre skozi shemo. Ta preveri, ali ima odgovor vsa
pričakovana polja in ali so ime zmote, kategorija zmote, vrsta zavrnitve in vrsta
izmikanja ena od dovoljenih vrednosti. Napačne vrednosti koda ne zavrže, ampak jo preslika na dovoljeno sopomenko. Tako na primer \texttt{false cause post hoc}
postane \texttt{post\_hoc}, \texttt{weak} postane \texttt{weak\_reasoning}, 
\texttt{undermining the premise} pa \texttt{undermining\_premise}. Če kakšno
polje manjka, dobi privzeto vrednost.

\label{ssec:arg-id}

Vsak argument dobi oznako iz imena govorca in zaporedne številke, na primer
\texttt{Speaker~1\#3}. Oznake dodeli sistem takoj po prvem koraku, ne model. Pri
istem izhodu prvega koraka so oznake zato vedno enake. 

Zmote in zavrnitve iz poznejših korakov se na argument sklicujejo z njegovo oznako.

\label{ssec:povezovanje}
Koda po koncu analize vsako oznako preveri. Pri zmoti mora oznaka pripadati argumentu istega govorca, pri zavrnitvi pa argumentu govorca, ki mu je zavrnitev namenjena. Če oznaka ne obstaja, koda izbere argument tega govorca, ki ima z navedkom zmote oziroma z opisom napadenega argumenta največ skupnih besed. Če takega argumenta ni, ostane oznaka prazna.

Kadar je odgovor modela odrezan, ga koda zahteva znova z dvakrat večjo omejitvijo izhodnih žetonov, kadar ni veljaven JSON, pa z enako, v obeh primerih skupaj največ trikrat. Če ne uspe prvi korak, se analiza ustavi, pri ostalih korakih pa se nadaljuje brez izida neuspelega koraka. 

\section{Stroški in ponovljivost}
\label{sec:stroski}

Obseg obdelave za posnetek dolžine 23 minut, kar je zgornja meja sistema,
povzema tabela~\ref{tbl:obseg}.

\begin{table}[H]
\centering
\small
\arrayrulecolor{crta}
\begin{tabular}{@{}p{0.62\textwidth}r@{}}
\toprule
Količina (posnetek 23 minut) & Vrednost \\
\midrule
izrečenih besed (skupaj vsi govorci) & \textcolor{red}{XX} \\
besed na minuto & \textcolor{red}{XX} \\
znakov prepisa & \textcolor{red}{XX} \\
žetonov prepisa & \textcolor{red}{XX} \\
\midrule
izluščenih argumentov & \textcolor{red}{XX} \\
premis & \textcolor{red}{XX} \\
zaznanih zmot & \textcolor{red}{XX} \\
preverjenih trditev & \textcolor{red}{XX} \\
virov skupaj (največ deset na trditev) & \textcolor{red}{XX} \\
\midrule
klicev jezikovnih modelov & \textcolor{red}{XX} \\
vhodnih žetonov skupaj & \textcolor{red}{XX} \\
izhodnih žetonov skupaj & \textcolor{red}{XX} \\
čas analize (min) & \textcolor{red}{XX} \\
strošek (EUR) & \textcolor{red}{XX} \\
\bottomrule
\end{tabular}
\caption{Obseg obdelave za posnetek dolžine 23 minut.}
\label{tbl:obseg}
\end{table}

Analiza takega posnetka porabi povprečno \textcolor{red}{XX} žetonov, traja \textcolor{red}{XX} minut in stane \textcolor{red}{XX}~EUR. Pri \textcolor{red}{XX} besedah na minuto nastane prepis s približno
\textcolor{red}{XX} žetoni, kar je \textcolor{red}{XX}~\% zgornje meje, ki jo
sistem dovoli za prepis. Največ žetonov porabi \textcolor{red}{XX}, največ časa
pa \textcolor{red}{XX}.


\section{Spletna aplikacija}
\label{sec:aplikacija}

\subsection{Vmesnik REST in izvajanje opravil v ozadju}
\label{ssec:backend}

Zaledni del je izveden v ogrodju FastAPI.
Analiza traja od nekaj minut do več deset minut, zato je ni mogoče izvesti znotraj enega zahtevka.
Oddaja posnetka zato ustvari opravilo, ki se izvede v svoji niti, odjemalec pa stanje spremlja s periodičnim povpraševanjem.

Vmesnik REST ima te končne točke:
\begin{itemize}
\item \texttt{POST /analyze} in \texttt{POST /analyze/upload}: asinhrona zahteva za analizo povezave oziroma naložene datoteke, ki vrne oznako opravila,
\item \texttt{GET /jobs/\{id\}}: trenutno stanje opravila in korak, ki teče,
\item \texttt{GET /debates/\{id\}}: shranjena analiza,
\item \texttt{PATCH /debates/\{id\}}: ročni popravki analize kot seznam operacij,
\item \texttt{POST /debates/\{id\}/rerun}: ponovna analiza nad shranjenim prepisom,
\item \texttt{POST /debates/\{id\}/recheck}: ponovno preverjanje dejstev nad shranjeno analizo,
\item \texttt{GET /debates/\{id\}/pdf}: poročilo v obliki PDF.
\end{itemize}

Vsako opravilo ima svojo delovno mapo. Ločena nit vsakih pet minut počisti opravila, ki so končana več kot eno uro: iz njihovih map izbriše zvok in vmesne datoteke, prepis pa obdrži. Ob zagonu strežnika enako počisti mape, ki jih nihče ni spremenil vsaj šest ur, tako da ne poseže v analizo, ki še teče.

Med izvajanjem odjemalec prikazuje, kateri korak teče, kar prikazuje
slika~\ref{fig:loading}. Splošen pogled na končano analizo prikazuje
slika~\ref{fig:izdelek}.

\subsection{Prikaz rezultatov}
\label{ssec:frontend}
Uporabniški vmesnik je enostranska aplikacija, izdelana z ogrodjem React.
Rezultat analize se prikaže kot časovnica po govorcih: vsak argument je
razčlenjen na premise in sklep, nanj vezana zavrnitev, zaznane zmote in
povezane preverjene trditve pa so prikazani ob njem. Vsak argument je na časovnici
oštevilčeno vozlišče, katerega oranžno polnilo pove, da je bila v argumentu
zaznana zmota ali neresnična trditev.

Na vrhu časovnice je povzetek sinteze, pod imenom vsakega govorca pa stavka o načinu argumentiranja in o razsodbah njegovih trditev. Opis vloge moderatorja koda združi z zapisom o moderatorju iz prvega koraka in ga prikaže v okvirju nad časovnico. Pri nastopu enega govorca so nepodprte trditve prikazane pod argumentom, s katerim imajo največ skupnih besed. Povzetek je tudi na začetku poročila PDF.

Nad časovnico so metapodatki analize, torej način, govorci, datum in povezava
do posnetka, ter pet dejanj: ročno urejanje, ponovna analiza nad
shranjenim prepisom, ponovno preverjanje dejstev nad isto analizo, izvoz v PDF
in pojasnilo, kako brati analizo. Prikazuje jih slika~\ref{fig:nastavitve}.
Zavihka \sn{Preverjanje dejstev} in \sn{Poročilo} prikazujeta seštevek razsodb in seznam vseh preverjenih trditev.

\begin{figure}[htbp]
\centering
\includegraphics[width=0.95\textwidth,keepaspectratio]{nastavitve.png}
\caption{Pet možnosti nad končano analizo, in sicer ročno urejanje, ponovna
analiza, ponovno preverjanje dejstev, izvoz v PDF in priročnik za uporabo.}
\label{fig:nastavitve}
\end{figure}

Zaznane zmote so naštete pod argumentom, v katerem so bile najdene, vsaka s
svojim imenom, kategorijo in delom argumenta, na katerem stoji. Bralec lahko
vsako zaznavo potrdi, zavrne ali odstrani iz analize. S tem je izpolnjena zahteva po nevtralnosti iz razdelka~\ref{sec:zahteve}.

Preverjene trditve so pripete k premisi, iz katere so nastale. Ob premisi stoji
pika v barvi razsodbe, klik nanjo pa odpre okno s trditvijo, razsodbo,
obrazložitvijo, seštevkom po petih razsodbah in oštevilčenim seznamom virov, v
katerem ima vsak vir svojo oznako. Oznaka vira ne ocenjuje vira samega, ampak
pove le, kam ta vir kaže glede preverjane trditve.

Prikaz argumentov po govorcih prikazuje slika~\ref{fig:vmesnik-arg1}, preverjene trditve z viri pa slika~\ref{fig:viri-okno}.

\subsection{Hramba, ponovna analiza in jeziki}
\label{ssec:rerun}

Prepisi se hranijo ločeno od delovnih map opravil, zato je mogoče analizo ponoviti brez ponovnega prenosa in ponovnega prepisa posnetka.
Ta možnost je bila ustvarjena za primere, ko se spremeni konfiguracija ali kateri od pozivov, izkazala pa se je za bistveno tudi pri ovrednotenju: brez nje bi bilo petkratno ponavljanje analize istega posnetka v razdelku~\ref{sec:ponovljivost} dražje.

\label{ssec:baza}
Za hrambo je uporabljen SQLite. Baza ima tri tabele: uporabnike, prijavne seje in
analize. Pri vsaki analizi se poleg metapodatkov hranita celoten izhod analize in
preverjanja dejstev v zapisu JSON. Prepisa baza ne hrani. Shranjen je v ločeni
mapi na disku, od koder ga uporabi ponovna analiza.

\label{ssec:i18n-pdf}

Vmesnik, poročilo in izvoz v PDF so na voljo v slovenščini in angleščini. Jezik se izbere ob zagonu analize. Pri slovenščini je vsem pozivom dodano navodilo, naj model besedilo v izhodu zapiše v slovenščini, vrednosti iz zaprtih množic pa ostanejo angleške, ker jih sistem obravnava v tem jeziku.

\subsection{Ročno urejanje analize}
\label{ssec:urejanje}

Ker se model lahko zmoti, lahko uporabnik analizo ročno popravi. Vmesnik podpira preimenovanje govorca, urejanje njegovega stališča, dodajanje, urejanje, premikanje med govorci in brisanje argumentov, dodajanje, urejanje in brisanje odgovorov ter zaznanih zmot, urejanje povzetka in naslova ter presojo zaznane zmote.
Pri vsaki zaznani zmoti bralec lahko označi, ali jo potrjuje ali zavrača, ne da bi jo pri tem izbrisal.

Vmesnik za urejanje prikazuje slika~\ref{fig:urejanje}.

\begin{figure}[H]
\centering
\includegraphics[width=\textwidth,height=0.95\textheight,keepaspectratio]{edit.png}
\caption{Vmesnik za urejanje analize. Zgoraj je naslov strani z gumboma Shrani in zapri. Spodaj je razdelek govorca s povzetkom zagovarjanega stališča. Pod njim je seznam argumentov. Vsak argument vsebuje: oznako govorca, besedilo sklepa, seznam premis z možnostjo urejanja, brisanja in dodajanja. Na koncu je razdelek za zavrnitve.}
\label{fig:urejanje}
\end{figure}

%%%%%%%%%%%%%%%%%%%%%%%%%%%%%%%%%%%%%%%%%%%%%%%%%%%%%%%%%%%%%%%%%%%%%%%%%%%%%%%
%  5  vrednotenje
%%%%%%%%%%%%%%%%%%%%%%%%%%%%%%%%%%%%%%%%%%%%%%%%%%%%%%%%%%%%%%%%%%%%%%%%%%%%%%%

\chapter{Vrednotenje}
\label{ch:vrednotenje}

Vrednotenje nam pove kako zanesljivo se je na izid mogoče opreti. Merjena je ponovljivost, torej kako podobni so si rezultati, ko sistem na istem gradivu teče petkrat. Večje kot je ujemanje, zanesljivejši je sistem.

\section{Zbirka in metodologija}
\label{sec:zbirka}

Meritve so bile opravljene na dveh posnetkih, eden v načinu solo, drugi v načinu razprave. Prvi je nastop ene govorke o tehnologiji v šolah, drugi pa razprava
dveh zgodovinarjev o antični Grčiji. Oba sta v angleščini in oba trajata manj kot 23 minut, kolikor je zgornja meja sistema.
% TODO: navedi natančni dolžini obeh posnetkov.
Izbrana sta bila zato, ker je pri obeh tema jasna, veliko je
sklicevanja na vire in statistične podatke, govorci se ne prekinjajo, posnetka pa sta kakovostna.

Vsak posnetek je analiziran petkrat, skupaj torej deset analiz. Prvi zagon vsake
serije opravi celoten cevovod, štirje nadaljnji pa uporabijo shranjeni prepis
(razdelek~\ref{ssec:rerun}). Prepis je konstanta in ni predmet meritve. Nihanje,
ki ga meritve zajamejo, izhaja iz analize. To je namerna omejitev, na katero se
vrača razdelek~\ref{sec:razprava}.

Meritve o izluščanju, temah in zmotah izhajajo iz teh desetih zagonov
neposredno. Preverjanje dejstev se je po njih spremenilo, zato je bilo nad istimi
desetimi shranjenimi analizami pognano znova. Argumenti, premise in zmote pri tem
ostanejo take, kot so bile izluščene, znova nastanejo samo trditve, viri in
razsodbe. Meritve v razdelku~\ref{sec:dejstva-meritev} izhajajo iz tega zagona.

Nastavitve so za vseh deset zagonov enake in navedene v
poglavju~\ref{ch:implementacija}.

Kot mera razpršenosti sta navedena vzorčni standardni odklon in koeficient
variacije, torej odklon glede na povprečje.

Vse, kar je v tem poglavju izmerjeno, je ponovljivost in ne pravilnost. Meritve
povedo, koliko se izidi sistema med zagoni ujemajo med sabo, ne pa, ali so
argumenti pravilno izluščeni, zmote pravilno poimenovane in razsodbe pravilno
izrečene. Sistem, ki bi se vsakič enako motil, bi po teh merilih dosegel
najboljši izid. Pravilnost bi zahtevala referenčno označeno zbirko, ki je za to
nalogo ni, na kar se vrača razdelek~\ref{sec:razprava}.

\section{Ponovljivost}
\label{sec:ponovljivost}

Ponovljivost pove, kako podobni so si rezultati analiz istega posnetka z istimi nastavitvami. Bolj kot so si analize med sabo podobne, bolj smo lahko prepričani, da sistem deluje konsistentno. Primerjamo število izluščenih argumentov, zaznane zmote, razsodbe dejstvenih trditev in to, ali se v vseh analizah pojavijo iste vsebinske teme.

\subsection{Število argumentov in premis}

Osnovni izid analize je množica argumentov. Tabela~\ref{tbl:ponovljivost}
prikazuje, kako se ta množica spreminja med zagoni nad istim prepisom.

\begin{table}[htbp]
\centering
\begin{tabular}{lrrrrr}
\toprule
& \multicolumn{5}{c}{zagon} \\
\cmidrule(l){2-6}
& 1 & 2 & 3 & 4 & 5 \\
\midrule
\multicolumn{6}{l}{\textbf{razprava dveh govorcev}} \\
argumentov          & 19 & 18 & 15 & 18 & 12 \\
premis              & 53 & 45 & 39 & 47 & 41 \\
premis na argument  & 2,8 & 2,5 & 2,6 & 2,6 & 3,4 \\
\midrule
\multicolumn{6}{l}{\textbf{nastop enega govorca}} \\
argumentov          & 3 & 3 & 3 & 4 & 3 \\
premis              & 19 & 19 & 17 & 22 & 16 \\
premis na argument  & 6,3 & 6,3 & 5,7 & 5,5 & 5,3 \\
\bottomrule
\end{tabular}
\caption{Število izluščenih argumentov in premis v petih zagonih nad istim
prepisom.}
\label{tbl:ponovljivost}
\end{table}

Pri razpravi se število argumentov giblje med 12 in 19, povprečje je 16,4 z
odklonom 2,9, koeficient variacije torej 17,6 odstotka. Število premis niha
manj, povprečno 45,0 z odklonom 5,5 oziroma 12,2 odstotka.

Pri nastopu enega govorca je nihanje v relativnem merilu podobno, in sicer 3,2
argumenta z odklonom 0,4, torej 14,0 odstotka, in 18,6 premise z odklonom 2,3,
torej 12,4 odstotka.

Med načinoma je razlika. Razprava dveh govorcev da petkrat toliko argumentov kot
enako dolg nastop enega, nastop pa dvakrat toliko premis na argument. Razprava
teče od teme do teme in vsak govorec zagovarja več ločenih tez, medtem ko
govorka v nastopu razvija majhno število tez z dolgim nizom razlogov. Sistem to
razliko izrazi, ne da bi bil za katero od zvrsti posebej nastavljen.

Nihanje števila argumentov pa ni nihanje vsebine. To pokaže naslednji razdelek.

\subsection{Razlika med zagoni}
\label{sec:ujemanje}

Pomembnejše od števila argumentov je, ali gre v vseh zagonih za isto vsebino. Argument je lahko drugače ubeseden, tema, ki jo obravnava, pa mora biti
ista.

Za razpravo je bilo opredeljenih deset tem, ki se v posnetku obravnavajo, za nastop pa štiri, in za vsak zagon ročno preverjeno, ali se tema pojavi med argumenti ali njihovimi premisami.
Rezultat je v tabeli~\ref{tbl:ujemanje}.

\begin{table}[H]
\centering
\begin{tabular}{lccccc c}
\toprule
tema & 1 & 2 & 3 & 4 & 5 & pojavitev \\
\midrule
\multicolumn{7}{l}{\textbf{razprava dveh govorcev}} \\
trojanska vojna              & X & X & X & X & X & 5/5 \\
Šparta in njen sloves        & X & X & X & X & X & 5/5 \\
atenska demokracija          & X & X & X & X & X & 5/5 \\
Ahil in Herkul               & X & X & X & X & X & 5/5 \\
Zevs                         & X & X & X & X & X & 5/5 \\
Odisej                       & X & X & X & X & ne & 4/5 \\
Aleksander Veliki            & X & X & X & X & X & 5/5 \\
Perzijci in Grki             & X & X & X & X & X & 5/5 \\
rimski uspeh                 & X & X & X & X & X & 5/5 \\
barve na kipih               & X & X & X & X & X & 5/5 \\
\midrule
\multicolumn{7}{l}{\textbf{nastop enega govorca}} \\
tehnologija v razredu        & X & X & X & X & X & 5/5 \\
telefoni in družbena omrežja & X & X & X & X & X & 5/5 \\
umetna inteligenca in nadzor & X & X & X & X & X & 5/5 \\
branje in pomnjenje          & X & X & X & X & X & 5/5 \\
\bottomrule
\end{tabular}
\caption{Pojavljanje vsebinskih tem v petih zagonih. Trinajst tem od
štirinajstih se pojavi v vseh petih zagonih svoje serije.}
\label{tbl:ujemanje}
\end{table}

Trinajst tem od štirinajstih se pojavi v vseh petih zagonih. Edina izjema je
Odisej, ki v petem zagonu razprave izpade.

Iz tega sledi, da razlika med 12 in 19 argumenti ni razlika v tem, o čem sistem poroča, ampak v
tem, kako drobno isto vsebino razdeli. Nihanje je v razmejitvi in ne v zaznavi.

\subsection{Zaznane zmote}
\label{sec:zmote-meritev}

Zmote so najmanj stabilen izid analize. Pri razpravi jih sistem zazna med dvema
in štirimi, povprečno 3,0 z odklonom 1,0, torej koeficient variacije 33,3
odstotka. Pri nastopu enega govorca med petimi in sedmimi, povprečno 5,6 z
odklonom 0,9, torej 16,0 odstotka. Podrobnosti povzema
tabela~\ref{tbl:zmote-stabilnost}.

\begin{table}[htbp]
\centering
\begin{tabular}{lrr}
\toprule
& razprava & nastop \\
\midrule
zaznav skupaj (5 zagonov)     & 15 & 28 \\
različnih imen                & 7 & 9 \\
kategorija šibko sklepanje    & 12 & 17 \\
kategorija neformalna         & 3 & 11 \\
kategorija formalna           & 0 & 0 \\
\bottomrule
\end{tabular}
\caption{Zaznane zmote v petih zagonih vsake serije. Nobena zaznava ne pade v
kategorijo formalnih zmot.}
\label{tbl:zmote-stabilnost}
\end{table}

Formalne zmote se ne pojavijo. Zaprti slovar vsebuje imena formalnih zmot in
koda kategorijo izpelje iz imena (razdelek~\ref{ssec:prehod2}), a v 43 zaznavah
ni nobene formalne. To ni nujno napaka sistema, saj je formalna zmota napaka v
obliki sklepanja in je v govorjeni razpravi redka, ker nihče ne govori v
silogizmih.

Najpogostejši imeni sta pri razpravi \sn{non sequitur} s sedmimi zaznavami od
petnajstih in pri nastopu \sn{post hoc} s sedmimi od osemindvajsetih. Obe sodita
v šibko sklepanje.

\section{Preverjanje dejstev}
\label{sec:dejstva-meritev}

\subsection{Število trditev}
\label{ssec:koliko-trditev}

Pri razpravi je bilo v petih zagonih preverjenih 298 trditev, na zagon
povprečno 59,6 z odklonom 12,3, torej koeficient variacije 20,6 odstotka. Pri
nastopu enega govorca 189 trditev, na zagon povprečno 37,8 z odklonom 2,4, torej
6,3 odstotka.

Načina se pri tem obnašata različno. Pri razpravi je nihanje števila trditev
večje od nihanja premis iz razdelka~\ref{sec:ponovljivost}, pri nastopu pa
občutno manjše. Pri razpravi se namreč sestavita dve presoji, saj se med zagoni
razlikuje najprej množica premis, nato pa še izbira, katere od njih so
preverljive. Pri nastopu je preverljivih podatkov manj in so bolj izraziti,
zato jih zagon zajame skoraj enako ne glede na to, kako so premise razdeljene.

\subsection{Porazdelitev razsodb}
\label{ssec:razsodbe-meritev}

Vseh preverjenih trditev je bilo 487. Pri 27 od njih razsodba ni uspela, kar
obravnava razdelek~\ref{ssec:neuspele}, zato je razsodbo dobilo 460 trditev,
282 pri razpravi in 178 pri nastopu. Deleži v tabeli so računani na teh 460 in
ne na 487, ker trditev brez razsodbe ne sodi v nobeno od petih vrednosti.

\begin{table}[H]
\centering
\begin{tabular}{lrrrr}
\toprule
& \multicolumn{2}{c}{razprava} & \multicolumn{2}{c}{nastop} \\
\cmidrule(lr){2-3}\cmidrule(l){4-5}
razsodba & trditev & delež & trditev & delež \\
\midrule
resnično           & 149 & 52,8 \% & 49 & 27,5 \% \\
delno resnično     & 92 & 32,6 \% & 78 & 43,8 \% \\
zavajajoče     & 21 & 7,4 \% & 37 & 20,8 \% \\
neresnično        & 14 & 5,0 \% & 7 & 3,9 \% \\
nepreverljivo  & 6 & 2,1 \% & 7 & 3,9 \% \\
\midrule
skupaj         & 282 & 100 \% & 178 & 100 \% \\
\bottomrule
\end{tabular}
\caption{Porazdelitev razsodb čez pet zagonov vsake serije. Deleži so računani
brez razsodb, ki niso uspele, te obravnava
razdelek~\ref{ssec:neuspele}.}
\label{tbl:zbirka}
\end{table}

Porazdelitvi se med posnetkoma razlikujeta in razlika je vsebinska. Razprava
zgodovinarjev navaja večinoma dobro dokumentirana dejstva o antiki, zato
prevladuje \sn{resnično} z 52,8 odstotka. Nastop o tehnologiji v šolah navaja
sodobne statistične in zdravstvene podatke, ki so pogosto izrečeni brez omejitev,
ki bi jih raziskava zahtevala, zato tam prevladujeta \sn{delno resnično} in
\sn{zavajajoče}, ki skupaj pokrivata 64,6 odstotka trditev.

Razsodba \sn{neresnično} je v obeh serijah redka, pri razpravi 5,0 in pri nastopu
3,9 odstotka. Trditev, izrečena v govoru, redko vsebuje vse omejitve, ki bi jo
naredile natanko točno, in redko je popolnoma napačna.

\subsection{Viri in njihove oznake}
\label{ssec:viri-meritev}

Vsaka od 460 trditev z razsodbo je dobila natanko deset virov, kolikor je
zgornja meja seznama (razdelek~\ref{ssec:viri}).  Število
različnih domen, ki jih teh deset virov pokriva je povprečno 8,0 pri
razpravi in 7,6 pri nastopu.

Model za razsodbo vsak vir posebej označi z eno od istih petih razsodb
(razdelek~\ref{ssec:razsojanje}). Porazdelitev teh oznak prikazuje
tabela~\ref{tbl:oznake-virov}.

\begin{table}[H]
\centering
\begin{tabular}{lrrrr}
\toprule
& \multicolumn{2}{c}{razprava} & \multicolumn{2}{c}{nastop} \\
\cmidrule(lr){2-3}\cmidrule(l){4-5}
oznaka vira & virov & delež & virov & delež \\
\midrule
nepreverljivo & 1790 & 63,5 \% & 1081 & 60,7 \% \\
resnično          & 634 & 22,5 \% & 305 & 17,1 \% \\
delno resnično    & 318 & 11,3 \% & 301 & 16,9 \% \\
neresnično       & 55 & 2,0 \% & 61 & 3,4 \% \\
zavajajoče    & 23 & 0,8 \% & 32 & 1,8 \% \\
\midrule
skupaj        & 2820 & 100 \% & 1780 & 100 \% \\
\bottomrule
\end{tabular}
\caption{Oznake posameznih virov čez pet zagonov vsake serije.}
\label{tbl:oznake-virov}
\end{table}

Skoraj dve tretjini virov o trditvi ne povesta ničesar. Povprečna trditev ima
med desetimi viri okoli štiri, ki se do nje vsebinsko opredelijo, pri razpravi
3,7 in pri nastopu 3,9. Iskalnik vrne strani, ki so tematsko blizu, kar pa še ne pomeni, da katera od
njih obravnava prav to število ali prav ta datum.

Ta porazdelitev tudi pokaže, zakaj seštevek oznak ni glasovanje
(razdelek~\ref{ssec:razsojanje}). Če se upoštevajo samo viri, ki se vsebinsko
opredelijo, se razsodba o trditvi v 67,8 odstotka primerov ujema z
najpogostejšo oznako med njimi, v 32,2 odstotka pa se od nje razlikuje. Skoraj
tretjina razsodb torej ni tista, ki bi jo dalo preštevanje virov. Razlika je
večja pri nastopu, kjer se razsodba od najpogostejše oznake razlikuje v 41,2
odstotka primerov, in manjša pri razpravi, kjer se v 26,5 odstotka. 

\subsection{Razsodbe, ki niso uspele}
\label{ssec:neuspele}

Kadar razsodbe ni mogoče dobiti, trditev ostane brez nje in se to zabeleži.
Takih trditev je bilo pri razpravi 16 od 298, torej 5,4 odstotka, in pri nastopu
11 od 189, torej 5,8 odstotka. Delež je v obeh serijah skoraj enak, kar pomeni,
da odpoved ni vezana na zvrst posnetka.

Vseh 27 neuspehov ima isti vzrok. Model vrne besedilo, ki ni veljaven JSON.

Ker razsodba, ki ni uspela, ni razsodba, so deleži v tabeli~\ref{tbl:zbirka} računani brez njih. Odpoved je vidna tudi uporabniku, saj trditev v poročilu
ostane brez razsodbe.

\section{Strošek in čas}
\label{sec:strosek-cas}

Čas in strošek obdelave povzema tabela~\ref{tbl:strosek}.

\begin{table}[htbp]
\centering
\begin{tabular}{lrr}
\toprule
& razprava & nastop \\
\midrule
\multicolumn{3}{l}{\textbf{čas obdelave}} \\
prvi zagon, celoten cevovod (min)   & 31,3 & 19,2 \\
ponovna analiza, povprečje (min)    & 11,9 & 8,2 \\
razpon ponovnih analiz (min)        & od 7,1 do 15,1 & od 7,0 do 9,6 \\
celotna serija petih zagonov (min)  & 79,0 & 52,1 \\
\midrule
\multicolumn{3}{l}{\textbf{strošek}} \\
prvi zagon, celoten cevovod (EUR)   & 1,78 & 1,16 \\
ponovna analiza, povprečje (EUR)    & 1,09 & 0,75 \\
prepis (EUR)                        & 0,69 & 0,41 \\
celotna serija petih zagonov (EUR)  & 6,14 & 4,16 \\
\bottomrule
\end{tabular}
\caption{Čas in strošek obdelave.}
\label{tbl:strosek}
\end{table}

Prvi zagon je pri razpravi trajal 31,3 minute, ponovne analize pa povprečno 11,9
minute. Razliko naredi prepis z ločevanjem govorcev, ki je časovno najdražja
stopnja cevovoda. Mehanizem ponovne analize (razdelek~\ref{ssec:rerun}) je s tem
pogoj za to vrednotenje, saj bi serija petih zagonov brez njega prepis plačala
petkrat.

Nihanje časa ponovnih analiz je precejšnje, od sedmih do 15 minut, in izhaja
pretežno iz zunanjih iskalnikov, ki niso pod nadzorom sistema.

Prva analiza je dražja od ponovnih, ker vključuje prepis.

\section{Razprava o omejitvah}
\label{sec:razprava}

Vse številke v tem poglavju merijo, ali se sistem ujema s samim sabo, ne pa,
ali je bistvo razprave korektno povzeto. Ali je posamezna zaznana zmota res zmota in ali je razsodba pravilna, ostaja stvar presoje bralca, ki ima za to v aplikaciji vgrajeno
orodje (razdelek~\ref{ssec:urejanje}).

Nihanje prepisa ni izmerjeno. Ker ponovne analize uporabljajo shranjeni prepis,
meritve zajamejo nihanje analize in ne nihanja prepisa.

Ena kategorija je prazna. Kategorija formalnih zmot se ni pojavila nikoli.
Zaprti slovar je torej širši od tega, kar sistem na tem gradivu dejansko
uporabi.

Del razsodb ne uspe. Odpoved razčlenjevanja odgovora zadene 27 trditev od 487,
torej 5,5 odstotka.

Vse meritve so bile opravljene na angleških posnetkih. Kako zanesljiva je
analiza slovenskih posnetkov, ni izmerjeno.

\section{Ugotovitve meritev}
\label{sec:sklepi-meritev}

Najbolj stabilna je izbira tem. Trinajst vsebinskih tem
od štirinajstih se pojavi v vseh petih zagonih svoje serije, kar pomeni, da
bralec dobi isti pregled predstavljenih tez ne glede na to, kateri zagon
prebere. Meritev pa pove le, da sistem vsakič najde iste teme, ne pa, da najde vse.

Manj zanesljiva je razmejitev teh tem na posamezne argumente, kjer nihanje znaša
17,6 odstotka. 

Najmanj zanesljiv je izid, ki zahteva presojo o kakovosti sklepanja. Zaznane
zmote nihajo pri razpravi za tretjino. Te izide je smiselno brati kot opozorila,
ki jih uporabnik preveri.

Pri preverjanju dejstev meritve pokažejo predvsem, da presoja ni preštevanje, saj
se skoraj tretjina razsodb razlikuje od najpogostejše oznake med viri, ki se do
trditve opredelijo. 

% 6  SKLEPNE UGOTOVITVE

\chapter{Sklepne ugotovitve}
\label{ch:sklep}

Izdelano orodje opravi osnovni pregled govora ali razprave. Izlušči argumente s premisami, označi logične zmote, poveže zavrnitve z napadenimi
argumenti in preveri dejstva v več virih. V ponovljenih
analizah istega posnetka so bili glavni sklopi argumentov prepoznani vsakič,
torej je pregled med zagoni stabilen. 

Človeška presoja pa ostaja potrebna. Ista teza je lahko zapisana v enem
argumentu ali v več, zato je razdelitev vedno tudi stvar razlage. Podobno velja
za zmote: ali je sklicevanje na avtoriteto zmota ali legitimen argument, je
odvisno od tega, kdo je avtoriteta in o čem govori. Sistem zato vsako zaznavo
opremi s tistim delom argumenta, v katerem naj bi bila napaka, in z razlago,
tako da jo uporabnik lahko potrdi ali zavrne. Aplikacija tudi ne razglasi
zmagovalca, ampak predstavi obe strani. Sodbo o razpravi sistem prepusti bralcu.

Največ dela je zahtevalo preverjanje, ali se koda in modeli ravnajo po zapisanih pravilih. Zahtevno je bilo tudi zapisati pravila za posamezni korak,
saj vnaprej ni bilo jasno, na kaj vse je treba model opozoriti.

Pri preverjanju dejstev sta za verodostojnost na voljo le število virov in število domen, ki pa ne povesta, kako trdna je razsodba. Če deset virov prihaja z malo znanih domen, je razsodba videti enako trdna kot tista, ki stoji na uradnih virih.

\section{Nadaljnje delo}
\label{sec:nadaljnje-delo}

Sistem pokriva široko področje, zato so možne izboljšave na vsakem koraku. Za nadaljnje delo predlagamo:

\begin{itemize}
\item \textbf{Daljši posnetki.} Sistem sprejme posnetke, krajše od 23 minut.
      Daljšo razpravo bi bilo treba razrezati na več delov in te nato bodisi
      znova združiti bodisi analizirati vsakega posebej, kar prinaša nove
      težave, na primer prevelik vhod za model ali ujemanje govorcev med deli.
\item \textbf{Več kot dva govorca.} Trenutno sta podprta nastop enega govorca
      in razprava dveh. Razprava dva proti dva bi zahtevala kompleksnejšo strukturo in drugačno preslikavo zavrnitev.
\item \textbf{Jeziki.} Vse meritve so bile opravljene na angleških posnetkih.
      Kakovost analize slovenskih in drugih posnetkov še ni izmerjena.
\item \textbf{Pravila glede na temo.} Vsi posnetki dobijo ista pravila.
      Zgodovinska in filozofska razprava imata drugačna težišča, zato bi
      prilagojena merila lahko izboljšala zaznavo zmot.
\item \textbf{Čas analize.} Analiza traja razmeroma dolgo, kar omejuje rabo,
      pri kateri bi si uporabnik želel izid takoj.
\end{itemize}

\cleardoublepage
\addcontentsline{toc}{chapter}{Literatura}
\bibliographystyle{plain}
\bibliography{literatura}

\appendix

\chapter{Modeli po korakih cevovoda}
\label{app:modeli}

Tabela~\ref{tbl:modeli} navaja, kateri jezikovni model opravlja katero nalogo v
cevovodu, in utemeljitev vsake izbire. Razporeditev je nastavljiva v datoteki
\texttt{config.yaml}, zato je učinek zamenjave posameznega modela mogoče
izmeriti ob nespremenjeni preostali kodi.

\begingroup
\small
\arrayrulecolor{crta}
\begin{longtable}{@{}p{0.28\textwidth}p{0.24\textwidth}p{0.40\textwidth}@{}}
\label{tbl:modeli}\\
\toprule
Korak & Model & Utemeljitev izbire \\
\midrule
\endfirsthead
\toprule
Korak & Model & Utemeljitev izbire \\
\midrule
\endhead
\bottomrule
\endfoot
prepis z ločevanjem govorcev & gpt-4o-transcribe-diarize & prepis in oznake govorcev v enem klicu, pri preizkusih natančnejši od drugih modelov \\
izluščanje argumentov (korak 1) & claude-sonnet-5 & razmejitev argumentov določa vse nadaljnje korake \\
zaznava zmot (korak 2) & claude-sonnet-5 & presoja odtenkov med legitimno potezo in zmoto \\
izluščanje preverljivih trditev & gpt-5.6-luna & visoke omejitve pretoka, zanesljiv strukturiran izhod \\
razgradnja sestavljenih trditev & gpt-4.1-nano & trivialna naloga, najcenejši razpoložljivi model \\
angleške ključne besede za akademske zbirke & gpt-4.1-nano & le pri trditvah v drugem jeziku, ker zbirke iščejo po angleških besedah \\
iskanje po spletu & gpt-4o & vgrajeno orodje za spletno iskanje \\
iskanje z navedbo virov & sonar-pro & iskanje po spletu z obveznim navajanjem virov \\
iskanje po spletu in omrežju X & grok-4.3 & vgrajeno iskanje po spletu in omrežju X \\
razsodba nad zbranim gradivom & claude-sonnet-5 & presoja gradiva brez lastnega iskanja \\
preslikava odgovorov (korak 4) & claude-haiku-4-5 & pretežno primerjalna naloga, nizka cena ob zadostni kakovosti \\
sinteza (korak 5) & claude-sonnet-5 & povzemanje čez izhode vseh korakov \\
\end{longtable}
\endgroup

\chapter{Ločnica med presojo in izračunom}
\label{app:locnica}

Tabela~\ref{tbl:locnica} po korakih razčleni, katere odločitve v cevovodu opravi
jezikovni model in katere koda. Odločitev je utemeljena v
razdelku~\ref{sec:kategorije-vs-stevilke}, izvedba pa je opisana v
razdelku~\ref{sec:determinizem}.

\begin{table}[htbp]
\centering
\small
\arrayrulecolor{crta}
\begin{tabular}{@{}p{0.44\textwidth} | p{0.44\textwidth}@{}}
\toprule
Model presodi & Sistem \\
\midrule
meje argumentov, stališča in premise & argumentom dodeli oznake \texttt{arg\_id}, odstrani argumente z manj kot dvema premisama \\
\midrule
ime zmote iz zaprtega slovarja & sopomenke poenoti, iz imena izpelje kategorijo \\
\midrule
del argumenta, v katerem je napaka, in njegovo oznako & preveri, da oznaka pripada istemu govorcu, sicer argument poišče po besedilu \\
\midrule
preverljive trditve in njihov tip & tip preslika na zaprto množico in po njem izbere zbiralce \\
\midrule
 & zbere gradivo, sestavi seznam največ desetih virov, prešteje vire in domene \\
\midrule
razsodbo nad gradivom in oznako vsakega vira & oznake preveri in sešteje \\
\midrule
vrsto zavrnitve in izmikanja & vrsti preslika na zaprto množico, zavrnitev pripne napadenemu argumentu \\
\midrule
 & zavrne posnetek, daljši od 23 minut, in pri razpravi prepis z enim samim govorcem, ustavi analizo, če pri razpravi ne ostaneta natanko dva debaterja \\
\bottomrule
\end{tabular}
\caption{Ločnica med presojo in izračunom.}
\label{tbl:locnica}
\end{table}

\chapter{Navodilo za izbiro preverljivih trditev}
\label{app:izbira-trditev}

Navodilo, po katerem model v tretjem koraku iz premis izbere preverljive trditve (razdelek~\ref{ssec:claim-extraction-fc}).

\begin{lstlisting}[basicstyle=\ttfamily\scriptsize,breaklines=true]
You are given the arguments already extracted from a debate, each with an
arg_id and numbered premises. Find the premises that assert something
CHECKABLE against outside sources and return them as claims.

A premise is checkable when it states a fact about the world: a number, a
date, a study, a measurable trend, who said or did what. A premise is NOT
checkable when it states what ought to be done, what is right or fair, or
how something feels - those are the positions being debated, not facts.

One premise may carry SEVERAL checkable facts (several statistics in one
sentence). Return each as its own claim, quoting that part of the premise as
written. Do not invent detail that is not in the premise, and do not
generalise it - carry over every number, date and scope word exactly as the
premise has them.

For each claim return:
- exact_claim: the checkable assertion, in the premise's own words
- arg_id: the arg_id it came from, EXACTLY as given
- premise_index: the number of the premise it came from
- speaker: the speaker of that argument
- claim_type: one of [statistic, historical, scientific, quote, policy,
  health, economic, geographic]
- context: one sentence on what the argument uses this fact for
\end{lstlisting}

\end{document}